\documentclass[11pt,oneside]{report}

% Thesis-length technical monograph assembled from the same source tree as
% paper/main.tex (the arXiv submission). This document is a synthesized
% internal artifact, not a work submitted to, or examined by, any degree-
% granting institution -- see the title-page note below. It reuses the
% paper's ggen/builder-rendered fragments by \input rather than restating
% any release standing number by hand, per this repository's own
% manufacturing law (AGENTS.md): nothing here is asserted in prose.

\usepackage[margin=1.15in]{geometry}
\usepackage[T1]{fontenc}
\usepackage[utf8]{inputenc}
\usepackage{amsmath,amssymb,amsthm}
\usepackage{booktabs}
\usepackage{microtype}
\usepackage[hidelinks]{hyperref}
\usepackage{url}
\usepackage{listings}
\usepackage{xcolor}
\usepackage{tocloft}
\usepackage{newunicodechar}

% The chapters quote Lean source fragments that use Unicode math notation
% directly (mathlib/procint convention); map the symbols that appear to
% their standard LaTeX math-mode equivalents so pdflatex can typeset them.
\newunicodechar{ℕ}{\ensuremath{\mathbb{N}}}
\newunicodechar{ℚ}{\ensuremath{\mathbb{Q}}}
\newunicodechar{≠}{\ensuremath{\neq}}
\newunicodechar{≤}{\ensuremath{\leq}}
\newunicodechar{⊕}{\ensuremath{\oplus}}

\lstdefinelanguage{lean}{
  keywords={structure, inductive, theorem, def, lemma, instance, namespace,
            end, import, where, by, fun, match, with, deriving},
  sensitive=true,
  comment=[l]{--},
  morestring=[b]",
}
\lstset{
  language=lean,
  basicstyle=\small\ttfamily,
  keywordstyle=\bfseries,
  commentstyle=\itshape\color{gray},
  columns=fullflexible,
  breaklines=true,
  frame=single,
  framerule=0.3pt,
  xleftmargin=1em,
}

\newtheorem{theorem}{Theorem}
\newtheorem{definition}{Definition}

% Release-mode guard, matching paper/main.tex's convention: release builds
% hard-fail on missing rendered fragments. There are no generated files --
% only artifacts with receipts: authority comes from .mfact/artifacts.toml,
% not from paths or headers. Volatile standing claims live only in
% ledgered, ggen-rendered fragments, \input exactly once here.
\newif\ifreleasebuild
\releasebuildtrue
% GENERATED FILE — DO NOT EDIT BY HAND
% Source: release-manifest.json + standing.env via the quadrature graph
% Generator: ggen (quadrature-pack) — ggen renders; Lean admits; mfact certifies.
% Standing: GENERATED_FROM_CERTIFIED_MANIFEST
\newcommand{\ReleaseTag}{v26.7.6-procint-certified}
\newcommand{\ReleaseHash}{a138ee84d0c08e0e946e0d0bb805a563b8304cf268eb97e6b9784bd36279fd86}
\newcommand{\ReleaseRun}{bdd8e99}
\newcommand{\SorryCount}{0}
\newcommand{\AdmitCount}{0}
\newcommand{\KernelTheorems}{145}
\newcommand{\TotalDecls}{318}
\newcommand{\StatedCount}{2}
\newcommand{\AxiomAuditStatus}{PASS}
\newcommand{\FixtureStatus}{PASS}
\newcommand{\CertifiedStatus}{PASS}
\newcommand{\QuadratureStatus}{PASS}
\newcommand{\CrownJewelStatus}{STATED}


\title{\Large Receipted Mathematical Manufacturing for Process Intelligence\\[0.3em]
\large A Thesis-Length Account of the mfact/procint Pipeline}
\author{Synthesized from the \texttt{mfact} repository\\(release tag \ReleaseTag)}
\date{\today}

\begin{document}

\maketitle

\thispagestyle{empty}
\vspace{-2em}
\begin{center}
\begin{minipage}{0.85\textwidth}
\small\itshape
\noindent Note on this document. This is a synthesized, thesis-length
technical monograph assembled from the \texttt{mfact} repository's own
source material (its LaTeX paper, doctrine documents, Lean declaration
catalog, and generated release fragments) together with agent-drafted
expository prose grounded in those sources. It has not been submitted to,
or examined by, any university or degree-granting body, and no academic
degree is claimed on its behalf. Its standing claims are exactly the
claims already certified by the repository's own manufacturing pipeline
(see Chapter~\ref{ch:mfact}); this document adds exposition, not standing.
\end{minipage}
\end{center}
\clearpage

\pagenumbering{roman}
\tableofcontents
\clearpage
\pagenumbering{arabic}

\chapter{Introduction}
\label{ch:intro}

\section{The Problem: Contact Is Not Consequence}

Formal verification has always separated the act of writing a proof from
the act of the kernel accepting it. What has changed in the last several
years is who, or what, is doing the writing. Large language models and
autonomous coding agents can now produce syntactically plausible Lean~4
declarations, tactic scripts, and entire modules on request. This
capability is genuinely useful, and the manufacturing system described in
this thesis was itself built with substantial LLM and agent assistance.
But it raises a question that the Lean kernel alone does not settle:
\emph{what standing does generated material have before, during, and
after it is admitted?}

Kernel acceptance is a strong guarantee of one specific thing: that a term
inhabits its stated type. It is not a guarantee that the stated type is
the intended one, that a multi-part statement was not silently weakened,
that a proof obligation was not narrowed by an extra hypothesis, or that
the artifact as a whole is a faithful rendering of some target
mathematical theory. Meek et al.~\cite{meek2026formalizing} catalog
exactly these failure modes in an agent-driven numerical-analysis
pipeline: compilation success can coexist with an added weakening
hypothesis, an incompletely stated multi-part theorem, or a parameter
restriction that trivializes the intended claim. None of these failures
are kernel failures. They are failures of the surrounding process --- the
part of the pipeline that decides which candidate gets typed, checked,
and recorded in the first place.

This repository's governing position is that an LLM, an agent, a
template engine, or a human editor is, with respect to the formal
corpus, an \emph{untrusted candidate producer}. AGENTS.md states this as
the repository's Prime Directive: ``This repository is a mathematical
manufacturing system. Do not hand-code manufactured outputs. Claude is an
untrusted candidate producer. Claude may propose changes to source
declarations, templates, fixtures, tests, and hand-authored prose. Claude
may not confer standing by directly editing generated artifacts.'' The
directive is deliberately mechanical rather than aspirational: it names
which files may be hand-edited (source declaration catalogs in Turtle,
projection templates, builder scripts, stable prose) and which may not
(rendered Lean modules, ledgered paper fragments, release manifests,
release hashes, and counts). The asymmetry is the point. A candidate
producer --- human or machine --- may contact the system at many points,
but contact does not by itself create a standing claim. Only a declared,
mechanically checked admission path can do that.

Put in the vocabulary this thesis adopts throughout: the standing law is
\[
R_B \;\vdash\; A = \mu(O^*_B),
\]
read inside a declared boundary~$B$: an artifact~$A$ with standing is the
lawful manufacturing transform~$\mu$ applied to admitted observations
$O^*_B$, and the receipt~$R_B$ is what makes that consequence checkable
rather than merely asserted. Raw candidate output --- a draft theorem
statement, a proposed tactic script, a hand-typed release count --- is
$O$, not $O^*$; it acquires no standing until it passes through the
admission gate that produces $O^*_B$. This separates \emph{contact}
(a candidate producer touching the corpus) from \emph{consequence} (a
kernel-checked, receipted, replayable change of standing). The remainder
of this thesis is an account of what it takes to make that separation
operational at the scale of a several-hundred-declaration formal
library, rather than as a slogan.

\section{Thesis Statement}

This thesis argues that formal artifacts produced with the assistance of
untrusted candidate producers can be given a disciplined, replayable
notion of \emph{standing} by treating formalization itself as a
manufacturing pipeline rather than as a sequence of ad hoc edits. The
pipeline has five stages, each of which produces or checks a distinct
kind of evidence rather than merely restating the previous stage's
output:

\begin{enumerate}
\item \textbf{Source declaration.} A domain's obligation graph --- the
  structures, theorem surfaces, module dependencies, fixtures, and
  citations a formal library must cover --- is curated as an RDF/Turtle
  catalog rather than written directly as Lean source. The catalog is the
  place where scope, naming, and status (``stated'' versus a status that
  a passing guard may promote toward ``proven'') are declared.
\item \textbf{ggen projection.} The declaration catalog is projected
  through SPARQL-in-Tera templates into candidate Lean~4 modules. This
  projection step is not itself an admission step: its output is a
  candidate, carrying provenance and a content hash, with no standing of
  its own.
\item \textbf{Lean kernel admission.} The projected modules are compiled
  under a pinned toolchain. Kernel acceptance, a \texttt{sorry} census,
  and an axiom audit together decide whether a candidate becomes an
  admitted, checked declaration --- and every failure mode is a typed
  refusal rather than a silently absorbed error.
\item \textbf{mfact certification.} A release manifest aggregates the
  admitted corpus's hashes, axiom footprints, fixture results, and
  proven/stated status into a single, machine-checkable object, closed
  against a declared vocabulary of admissible objections.
\item \textbf{Replayable receipt.} The manifest, the axiom audit log, and
  the negative fixtures together constitute a receipt: evidence that a
  claimed artifact is a genuine consequence of the pipeline, checkable
  again by a third party rather than trusted on the strength of the
  first run.
\end{enumerate}

This thesis further argues that this discipline is not confined to
abstract pipeline design: it is exercised on a substantial first
manufactured vertical, \textsf{procint}, a Lean~4 library for process
intelligence --- Petri nets, workflow-net soundness, event logs and
alignment-based conformance, and OCEL-style object-centric event
data --- manufactured from the declared obligation graph rather than
hand-written declaration by declaration. Process intelligence is chosen
not because it is the terminus of the method but because, as
paper/main.tex's Scope paragraph is explicit in stating, its public
objects (events, traces, transitions, guards, conformance obligations)
already carry mathematical structure suitable for this kind of projection
and admission, making it a demonstration vertical for a method the thesis
argues is field-independent in its structure, if not in this document's
empirical coverage.

\section{Contributions}

This thesis makes the following contributions, each developed in a later
chapter:

\begin{enumerate}
\item \textbf{A typed admission discipline for formal artifacts
  (\emph{mathematical manufacturing}).} Generated output, human patches,
  templates, and agent proposals are all treated uniformly as
  candidates; standing is acquired only through an explicit, fail-closed
  admission gate --- kernel acceptance, a zero-\texttt{sorry} census, an
  axiom audit against an authorized set, negative fixtures, and manifest
  completeness (Chapter~\ref{ch:background}, Chapter~\ref{ch:mfact}).
\item \textbf{\textsf{mfact}, a Lean~4 manufacturing framework} that
  reifies candidates, refusals, admissions, manifests, certified
  releases, and objections as first-class typed values rather than as
  informal process conventions, and that discharges a per-release
  ``no valid objection'' theorem against a closed inductive family of
  admissible objection forms (Chapter~\ref{ch:mfact}).
\item \textbf{\textsf{procint}, the first manufactured vertical}: a
  process-intelligence formalization covering Petri nets, workflow-net
  soundness (including the crown-jewel soundness equivalence, proven
  under a \texttt{[Finite T]} hypothesis, and a deliberately unbounded
  countermodel witnessing why that hypothesis cannot be dropped), event
  logs, alignment-based conformance, and OCEL~2.0 object-centric logs,
  together with an account of the exactly one declaration in the current
  \TotalDecls-declaration catalog that remains genuine open proof work
  (Chapter~\ref{ch:procint}).
\item \textbf{A correspondence rail from extracted implementation to
  admitted semantics}, using Aeneas/Charon to extract a bounded Rust core
  into a proof-facing image and relating that image to admitted
  \textsf{procint} semantics under a declared, kernel-checked abstraction
  function, rather than by prose or convention alone
  (Chapter~\ref{ch:correspondence}).
\item \textbf{A manifest-generated, kernel-closed evaluation}: the
  paper's own numeric claims are rendered from the release manifest, not
  asserted in prose, and closed by a kernel-checked Standing Quadrature
  witness that refuses the release if the declaration catalog, the
  admitted corpus, and the manifest ever diverge (Chapter~\ref{ch:empirical},
  Chapter~\ref{ch:evaluation}).
\item \textbf{An account of the pipeline's own construction} by agentic
  and LLM-assisted development, examined under the same untrusted-producer
  discipline the pipeline enforces on its mathematical content
  (Chapter~\ref{ch:ai-development}).
\end{enumerate}

\section{Scope and Non-Claims}

Consistent with the Scope paragraph of the accompanying paper, this
thesis does not claim to formalize process intelligence in its entirety,
nor does it claim the manufacturing pipeline's benefits transfer to every
formal domain without further work; the generalization argument
(Chapter~\ref{ch:discussion}) is presented as a structural analogy backed
by one worked vertical, not as an empirical result over multiple
verticals. It does not claim that generated mathematics is correct by
virtue of being generated, or that kernel admission settles the harder
semantic question of whether an admitted statement is a faithful
rendering of its informal target --- a residual risk this thesis returns
to explicitly, following Meek et al.'s diagnosis, in
Chapter~\ref{ch:discussion}. It does not claim that every declaration in
the corpus is proven: the release manifest itself distinguishes
kernel-admitted, axiom-audited theorems from declarations whose formal
statement exists but whose proof obligation remains open, and this
thesis is careful to preserve that distinction rather than to elide it in
narrative prose. Where a number is asserted --- a declaration count, a
theorem count, a hash --- it is drawn from a rendered, ledgered fragment
of the release manifest rather than typed by hand into this document.

\section{Roadmap}

The remaining chapters develop the contributions above in the following
order. Chapter~\ref{ch:background} situates this work against the Lean~4
ecosystem, Petri-net and workflow-net theory, the process-mining
literature, LLM-assisted theorem proving, and Rust verification and
extraction, closing with an explicit statement of the research gap this
thesis addresses. Chapter~\ref{ch:mfact} develops the \textsf{mfact}
framework itself: candidates, typed refusals, admissions with covers
proofs, manifests, certified releases, and the closed family of valid
objections. Chapter~\ref{ch:procint} presents \textsf{procint} as the
first manufactured vertical, including its module architecture, the
workflow-net soundness development, and the deliberate countermodel that
bounds the crown-jewel theorem's scope. Chapter~\ref{ch:correspondence}
develops the Aeneas/Charon correspondence rail between extracted Rust
implementation and admitted process semantics. Chapter~\ref{ch:empirical}
reports the manifest-generated evaluation and the Standing Quadrature
closure argument in detail, and Chapter~\ref{ch:ai-development} turns the
thesis's own untrusted-producer discipline onto the agentic process that
built the pipeline. Chapter~\ref{ch:evaluation} consolidates the
quantitative release evidence referenced throughout the thesis into a
single evaluation treatment, and Chapter~\ref{ch:discussion} discusses
limitations, the residual semantic-faithfulness risk, and the
generalization argument to ontology-bearing fields beyond process
intelligence. Chapter~\ref{ch:conclusion} closes the thesis.

\chapter{Background and Related Work}
\label{ch:background}

This chapter surveys the four bodies of work this thesis draws on and
positions itself against: the Lean~4 ecosystem and its formal-library
practice; Petri-net and workflow-net soundness theory together with the
process-mining literature it feeds; the growing landscape of
LLM-assisted and automated theorem proving; and Rust verification and
extraction. It closes with an explicit statement of the gap between
these tracks that motivates this thesis's contribution.

\section{Lean 4 and Its Ecosystem}

Lean~4 itself, as a dependently typed theorem prover and programming
language with an elaboration and tactic framework distinct from Lean~3,
is described by de Moura and Ullrich~\cite{moura2021lean4}. The design
choice most relevant to this thesis is that Lean~4's kernel is small and
independently re-implementable: Lean4Lean, an independent kernel-in-Lean
typechecker, has been used to re-verify the reference kernel's behavior
and, in doing so, located and helped fix a soundness bug in the
reference implementation~\cite{lean4lean2024}. This matters for the
manufacturing pipeline described in this thesis because the pipeline's
own trust boundary is explicitly drawn \emph{above} the kernel: the
pipeline does not attempt to re-verify or extend kernel soundness, it
relies on the kernel as an already-scrutinized primitive and confines its
own apparatus (candidates, refusals, admissions, manifests) to the layer
of process around kernel calls.

Mathlib~\cite{mathlib2020} is the reference point for what a
large-scale, collaboratively maintained Lean formal library looks like in
practice: its quality regime is human review plus continuous
integration, applied uniformly across a very large and heterogeneous
contributor base. \textsf{procint} depends on Mathlib for its background
mathematics (finite types, order structures, and the like), but does not
attempt to replicate Mathlib's review-based quality regime; instead it
substitutes a mechanically checked admission gate for human review as the
primary standing-conferring step, a substitution this thesis argues is
appropriate specifically because \textsf{procint}'s corpus is
manufactured from a declared, machine-readable obligation graph rather
than authored freely.

CSLib~\cite{cslib2026} extends the Lean ecosystem toward computer
science proper --- labeled transition systems, traces, and bisimulation
among its concerns --- and \textsf{procint} builds directly on this layer
for parts of its semantics, in particular where process behavior is
described via transition and trace structures rather than purely via
algebraic Petri-net state equations. The presence of a computer-science
library alongside Mathlib's predominantly mathematical corpus is itself
evidence that Lean's applicability is already understood to extend past
pure mathematics; \textsf{procint} extends that applicability further,
into a domain --- process intelligence --- that neither Mathlib nor
CSLib targets directly.

\section{Petri Nets, Workflow-Net Soundness, and Process Mining}

The structural theory of Petri nets that \textsf{procint}'s
\texttt{ProcInt.Petri} namespace formalizes --- the state equation,
place and transition invariants, boundedness, and liveness --- follows
Murata's survey~\cite{murata1989}, still the standard reference
articulation of these properties and their relationships (for instance,
that liveness implies deadlock-freedom but not conversely).

Workflow nets, and the soundness property this thesis's crown-jewel
theorem addresses, are due to van der Aalst~\cite{vanderaalst1997}. The
specific result \textsf{procint} formalizes is van der Aalst's Lemma~8:
for a workflow net with a single source place, a single sink place, and
every node on some path from source to sink, soundness is equivalent to
liveness and boundedness of the net's short-circuited variant (the net
obtained by adding a transition from the sink back to the source). This
equivalence is exactly the theorem this thesis refers to throughout as
the crown-jewel theorem; Chapter~\ref{ch:procint} gives its Lean
statement and admission status in detail. The subsequent decidability
and classification landscape for workflow-net soundness --- which
soundness variants are decidable, and under what net restrictions --- is
surveyed in~\cite{vanderaalst2011soundness}, and the broader process
mining setting these soundness notions serve (discovery, conformance,
enhancement) is presented as a textbook treatment
in~\cite{vanderaalst2011pm}.

\textsf{procint}'s conformance-checking layer, which relates a process
model to observed event-log behavior via alignments, moves, and cost,
follows Carmona et al.'s treatment of conformance
checking~\cite{carmona2018}, which frames alignment-based techniques as
the technically preferred (if computationally heavier) alternative to
naive token-replay methods. Object-centric event data --- where a single
event may relate to several objects of possibly different types, rather
than to a single case identifier --- follows the OCEL~2.0
specification~\cite{ocel2}, and \textsf{procint}'s \texttt{ProcInt.Ocel}
namespace formalizes the entity-to-entity and object-to-object relations,
flattening, and lifecycle notions that specification defines.
Object-centric Petri nets, the net-level counterpart to object-centric
event logs, follow~\cite{ocpn2020}. Partially ordered process models,
formalized alongside process trees, DFGs, causal nets, BPMN, and Declare
in \textsf{procint}'s \texttt{ProcInt.Models} namespace, follow the POWL
language of Kourani and van Zelst~\cite{powl2023}.

To the best of a bounded search of the Lean, Coq, and Isabelle/AFP
library indexes conducted for this project (a negative search result,
not a proof of absence), none of workflow-net soundness, alignment-based
conformance, or OCEL had a prior public mechanization in a proof
assistant at the time \textsf{procint} was manufactured. A separate,
internal literature scan conducted for this project and recorded in this
repository's own research notes (\texttt{research-papers/README.md})
reaches a related but distinct observation: across the papers that scan
surveyed, process mining and formal (theorem-prover) verification have
historically progressed along near-parallel tracks, with comparatively
few works that bridge the two --- the scan's own phrasing is that the two
areas ``operate in parallel'' with essentially no bridge papers located.
That internal document is a project research note rather than a
peer-reviewed source, and this thesis treats its finding accordingly, as
a motivating internal observation rather than as a citable external
claim; but it is consistent with, and independently arrived at from, the
absence this thesis reports from the Lean/Coq/Isabelle index search
above.

\section{LLM-Assisted and Automated Theorem Proving}

A substantial and fast-moving literature addresses how language models
and automated provers can assist, rather than replace, interactive
theorem proving in systems like Lean. Two architectural patterns
recur. The first is in-process integration, where a language model is
called directly from within the tactic framework via a foreign-function
boundary, as in Lean Copilot~\cite{leancopilot2024}. The second is
external bridging to dedicated automated theorem provers or solvers
outside the interactive session, as in Lean-auto~\cite{leanauto2025},
which interfaces Lean~4 with external ATPs. \textsf{mfact} adopts neither
pattern as a trust mechanism: whatever automation --- language-model or
otherwise --- is used to produce a candidate declaration or proof
script, the candidate is untrusted data with respect to the pipeline's
own trust boundary, admitted only through the gate developed in
Chapter~\ref{ch:mfact}. This is a deliberate design choice rather than an
oversight of the in-process/external-bridge distinction: from the
pipeline's point of view, an in-process copilot suggestion and an
external solver's output are the same kind of object, a candidate, and
neither is privileged relative to a template-rendered or a
hand-authored one.

Search-based and learning-based automated proving more broadly is
represented by HOList~\cite{holist2019}, an environment for machine
learning applied to higher-order theorem proving, and by HyperTree Proof
Search~\cite{hypertree2022}, which frames proof search itself as a
tree-search problem amenable to neural guidance. TheoremLlama~\cite{theoremllama2024}
and APOLLO~\cite{apollo2025} represent, respectively, adapting a
general-purpose LLM into a Lean~4-specialized prover and orchestrating
collaboration between an LLM and Lean's own checking loop toward more
advanced formal reasoning. Infrastructure supporting these efforts
includes premise and declaration search over existing libraries, as in
LeanExplore~\cite{leanexplore2025}, and large synthetic theorem corpora
generated by exploring state-transition graphs, as
in~\cite{mathlib4dataset2025}. \textsf{procint} does not currently have
the problems this infrastructure addresses: its statement corpus is
curated from a declared RDF/Turtle catalog rather than mined from an
existing library or synthesized at scale, so premise search over a large
undifferentiated corpus and synthetic-theorem generation are, for this
project's current scope, orthogonal concerns rather than adjacent
solutions.

Two further strands bear directly on this thesis's central claim that
kernel admission alone is an insufficient standard for LLM-produced
formal content. Programming Language Confusion~\cite{langconfusion2025}
documents that code-generating LLMs can fail to keep target languages
straight even when producing superficially well-formed output, an
instance of the broader concern that syntactic plausibility and semantic
fidelity can diverge under automated generation. Work on understanding
and improving automated proof synthesis for interactive theorem
provers~\cite{tacticexpertalign2026} studies exactly the alignment
between what a synthesized tactic sequence proves and what a human
intended it to prove. Most directly, Meek et al.~\cite{meek2026formalizing}
formalize numerical analysis using a coding-agent pipeline and audit its
output for quality beyond kernel acceptance, documenting concretely that
compilation success can mask an incompletely stated multi-part theorem,
a silently added weakening hypothesis, or an unannounced parameter
restriction. Meek et al. is the closest prior work to \textsf{mfact} in
problem framing, and this thesis differs from it in mechanism rather
than diagnosis, as elaborated in the Research Gap section below.

\section{Rust Verification and Extraction}

The correspondence rail developed in Chapter~\ref{ch:correspondence}
extracts a bounded Rust implementation core into a Lean-facing semantic
image via Aeneas~\cite{ho2022aeneas}, a tool for Rust verification by
functional translation, whose Charon front end~\cite{protzenko2024charon}
separates Rust syntax recovery from the semantic translation step
proper. Aeneas is designed to target whole Rust programs under a full
memory model; this thesis's use of it is narrower, applied to a single,
dependency-free arithmetic module, and --- consistent with the
untrusted-producer discipline argued for throughout this thesis ---
extraction's output is treated as a candidate Lean statement, never as a
proof: standing for a correspondence claim still requires kernel
admission of a hand-authored correspondence theorem stated over the
extracted declaration, not merely a same-shaped statement over
independently hand-written types. Chapter~\ref{ch:correspondence}
returns to this point in more mechanical detail, including the guard
that blocks a correspondence obligation from reporting proven status
while its extraction binding remains incomplete.

\section{Research Gap}

Taken individually, none of the tracks surveyed above addresses the
problem this thesis targets. The Lean ecosystem (Mathlib, CSLib,
Lean~4Lean) supplies library scale, computer-science-relevant semantics
infrastructure, and an independently checked kernel, but its standard
quality regime is human review, not a mechanically checked admission
pipeline applied uniformly to agent-, template-, and human-produced
candidates alike. The Petri-net, workflow-net, and process-mining
literature supplies the mathematical canon --- soundness, alignments,
OCEL --- that this thesis's first manufactured vertical formalizes, but
that literature is, per both this project's own bounded library-index
search and its internal research scan, essentially disjoint from formal
verification practice: process mining and theorem proving have
progressed as near-parallel tracks with few bridge works, so no existing
mechanization supplies workflow-net soundness or OCEL as a Lean~4 target
to build on or compare against directly. The LLM-assisted and automated
theorem-proving literature supplies both the capability this thesis takes
as its starting condition (agents and models can produce plausible formal
candidates) and, in Meek et al.~\cite{meek2026formalizing} particularly, a
sharp diagnosis of the resulting risk, but that diagnosis is delivered as
a post-hoc semantic audit layered over agent output after the fact, not
as a typed, structurally enforced pipeline that a candidate must pass
through before a standing claim can be recorded at all. The Rust
verification and extraction literature (Aeneas, Charon) supplies a
mechanism for turning implementation code into a proof-facing semantic
object, but does not by itself address how the resulting extracted
statement's standing, or the standing of any hand-authored correspondence
theorem built on it, should be tracked, gated, or receipted over time.

The gap this thesis addresses sits at the intersection of these tracks:
a domain formalization effort (process intelligence) that (a) has, to
this project's search, no prior proof-assistant mechanization to inherit
conventions from; (b) is produced with substantial LLM and agent
assistance, and therefore inherits the standing problem Meek et al.
diagnose; and (c) is manufactured, rather than hand-authored, from a
declared obligation graph, so that the standing problem must be solved
structurally --- as a typed admission pipeline with typed refusals,
manifests, and replayable receipts --- rather than remedied afterward by
a separate audit pass. Where Meek et al.'s audit is a quality assessment
performed \emph{over} already-produced agent output, this thesis's
contribution, \textsf{mfact}, is a pipeline in which candidates,
refusals, and admissions are themselves Lean data types, so that the
proven/stated distinction is enforced by a release gate rather than
estimated after the fact. The two approaches are, as
paper/main.tex's Related Work section observes, complementary rather
than competing: the residual risk that a kernel-admitted statement is
nonetheless a faithless or vacuous rendering of its informal target is
exactly the risk a semantic audit like Meek et al.'s is suited to catch,
and exactly the risk this thesis's pipeline, by design, does not claim
to eliminate. Chapter~\ref{ch:discussion} returns to this residual risk
directly, rather than treating kernel admission as a substitute for it.

\chapter{The mfact Manufacturing Framework}
\label{ch:mfact}
\label{sec:mfact}

Chapter~\ref{ch:background} situated this thesis against Lean's ecosystem
and the wider LLM-assisted-proving literature. This chapter turns to the
apparatus the thesis actually builds: \textsf{mfact}, the Lean~4 layer
that gives the manufacturing pipeline's epistemics --- candidate,
refusal, admission, manifest, certified release, valid objection --- a
first-class typed representation rather than leaving them as informal
convention. As paper/main.tex states it, \textsf{mfact} is deliberately
small: seven modules and a command-line interface. It is not the whole
system. It is the Lean/Lake standing apparatus that sits inside a larger
standing-manufacturing graph whose other rails (process evidence,
benchmark evidence, correspondence evidence, publication evidence) are
developed in later chapters. This chapter develops, in order, the
governing law the whole pipeline answers to; the typed vocabulary of
candidates and refusals; admission with a covers proof; the manifest,
certified release, and closed-objection theorem; the ledger law that
locates authority in a machine-checkable ledger rather than in file
paths or naming conventions; and the agent actuation constitution that
governs how any producer --- human or machine --- is permitted to touch
this apparatus.

\section{The Receipted Manufacturing Law}
\label{sec:mfact-receipted-law}

The governing law of the pipeline, stated once in
Chapter~\ref{ch:intro} and elaborated here, is
\[
A = \mu_B(O^*_B), \qquad R_B = \mathrm{receipt}(A).
\]
Every symbol names a distinct, checkable object rather than a stage of
prose. $B$ is the declared boundary: the Lean packages, the declaration
catalog, the manifest schema, and the audited theorem set that together
fix what counts as ``inside'' for a given release. $O$ is raw
observation --- candidate output from a human editor, a language model,
a ggen template, or a repository fragment, with no standing of its own.
$O^*_B$ is admitted observation inside $B$: typed, kernel-checked,
bounded, and replayable, the only class of observation the manufacturing
transform $\mu_B$ is permitted to consume. $\mu_B$ is the lawful
manufacturing transform itself --- ggen rendering, Lean kernel
admission, the Lake build, and the axiom/sorry audit, all parameterized
by the boundary $B$ they operate inside. $A$ is the resulting artifact
with standing: an admitted theorem, a module, a manifest entry, a
release. $R_B$ is the receipt --- the build log, the manifest hash, the
axiom audit output, the negative-fixture result, the replay evidence ---
that proves $A$ is a genuine consequence of $\mu_B$ acting on $O^*_B$,
rather than merely an assertion made about $A$ after the fact. The
asymmetry this equation encodes is the one already stated in
Chapter~\ref{ch:intro}: raw generated output is $O$, never $O^*$, until
it has passed through an admission step; manufacture produces the
artifact, the receipt proves the consequence relation held, and replay
proves the receipt was non-ornamental --- that a third party, not just
the original run, can recompute the same result. As paper/main.tex puts
it, the pipeline exists to refuse the gap between contact and
consequence.

Within this law the manufacturing graph distinguishes four kinds of
evidence, each answering to a different admission boundary rather than
being an interchangeable substitute for the others: \emph{formal}
evidence (Lean theorems and their axiom audits, the subject of this
chapter and Chapter~\ref{ch:procint}); \emph{operational} evidence (the
praxis-graphlaw benchmark and its execution validation via rslab,
Chapter~\ref{ch:empirical}); \emph{correspondence} evidence
(Aeneas-extracted implementation image related to admitted \textsf{procint}
semantics, Chapter~\ref{ch:correspondence}); and \emph{publication}
evidence (generated paper fragments and manifests, of which this
thesis's own numeric claims are themselves instances). The
\texttt{final\_status.tex} fragment below reports the current release's
standing across several of these evidence kinds at once, rendered
directly from the manifest rather than typed here.

\begin{table}[h]
\centering
\input{../paper/final_status}
\caption{Current release standing, rendered from the manifest via \texttt{final\_status.tex}.}
\label{tab:mfact-final-status}
\end{table}

In this architecture ggen is not a code generator in the ordinary sense.
It is the projection half of $\mu_B$: it maps a declared obligation
graph into candidate Lean surfaces. The Lean kernel supplies the
admission half. Projection without admission confers no standing;
admission without a projection mechanism does not scale to a
several-hundred-declaration corpus. Neither half is optional, and
neither half does the other's job --- a point returned to below when
Section~\ref{sec:mfact-ledger} distinguishes what a template may render
from what only kernel admission may certify.

\paragraph{Rice's boundary and manufactured decidability.} Rice's
theorem bars a nontrivial semantic decision procedure over arbitrary
programs, and the manufacturing law does not evade this boundary; it
changes the object of judgment. The pipeline does not decide the
semantic correctness of an arbitrary generated artifact considered in
isolation. It decides whether a bounded candidate has been admitted by a
specific Lean kernel, under a declared axiom policy, with a complete
manifest receipt --- a decidable question that replaces an undecidable
one rather than answering it. Undecidability is fenced, not defeated,
and the standing this apparatus manufactures is standing \emph{inside
that fence}, a scope this chapter and Chapter~\ref{ch:discussion} both
return to when discussing what the pipeline does not claim.

\paragraph{Status is algebra, not prose.} The release distinguishes a
closed vocabulary of standing values --- \textsc{Declared},
\textsc{Extracted}, \textsc{Stated}, \textsc{Proven}, \textsc{Refused},
\textsc{Blocked}, \textsc{BuildBroken}, \textsc{Unknown},
\textsc{Unsupported}, and \textsc{Alive} --- and no transition between
them is implicit or inferred from surrounding prose. \textsc{Extracted}
is not proof; \textsc{Unknown} is not admitted; \textsc{Unsupported} is
not the same refusal as \textsc{Refused}. A candidate artifact, whatever
its origin, remains exactly that --- a candidate --- until an explicit
admission step assigns it one of these values, and the repository's
broader status taxonomy (\textsc{Alive}, \textsc{Partial\_Alive},
\textsc{Blocked}, \textsc{Blocked\_External}, \textsc{Build\_Broken},
\textsc{Refused}, \textsc{Pending\_External\_Actuation}, \textsc{Valid},
\textsc{Unsupported}, \textsc{Planned}, \textsc{Stated},
\textsc{Replay\_Not\_Run}, \textsc{In\_Progress}, per AGENTS.md) refines
this same discipline at the level of whole packets and reports, not just
individual theorems.

\section{Candidates and the Typed Refusal Family}
\label{sec:mfact-candidates}

A \texttt{Candidate} in \textsf{mfact} is content plus provenance ---
producer identity and a content hash --- reflecting that every producer,
human or automated, is untrusted by construction with respect to the
formal corpus; nothing about a candidate's origin changes what it is
permitted to do. A failed admission is not represented as an untyped
error string but as a value of a closed inductive type,
\texttt{Refusal}, whose constructors --- \texttt{kernelRejected},
\texttt{sorryPresent}, \texttt{unauthorizedAxiom},
\texttt{missingEvidence}, \texttt{malformedModel},
\texttt{invalidFixture} --- name specific, checkable failure modes.
Because the type is closed and carries no default case, an unrecognized
failure is itself required to be classified as a refusal rather than
silently passed through; there is no code path by which a candidate can
fail to be admitted without producing a typed, loggable, re-queueable
record of why.

This same discipline is carried into the repository's engineering
practice at a coarser grain, as a vocabulary of typed refusals a human
or agent contributor invokes when a proposed change would confer
standing without a corresponding admission step:
\texttt{HAND\_CODED\_GENERATED\_OUTPUT} (a generated artifact was
hand-edited directly), \texttt{GENERATED\_OUTPUT\_DRIFT} (a rendered
file no longer matches its declared source), \texttt{MISSING\_GGEN\_SOURCE}
and \texttt{MISSING\_GGEN\_TEMPLATE} (an artifact is ledgered as
generated but its producing source or template cannot be found),
\texttt{ORPHAN\_GENERATED\_FILE} and \texttt{UNREGISTERED\_PAPER\_FRAGMENT}
(a file carries the shape of a generated or paper artifact but is absent
from the ledger), \texttt{UNSUPPORTED\_STANDING\_CLAIM} and
\texttt{STATED\_PROMOTED\_TO\_PROVEN} (a claim exceeds its evidence, or a
status transition happened without the checked guard that must license
it), \texttt{MANUAL\_RELEASE\_COUNT} and \texttt{MANUAL\_RELEASE\_HASH}
(a release number or hash was typed by hand instead of read from a
generated fragment), and \texttt{RECEIPT\_RECURSION\_REFUSED} (an
operation would mutate the identity of a release it is supposed to be
reporting on, discussed further in
Section~\ref{sec:mfact-core-identity}). A further code,
\texttt{SOURCE\_CHANGE\_ASSERTION\_UNSUPPORTED}, is planned to require
that any claim of a changed source ship with an actual changed
source, template, or TTL fragment in the same commit, not a bare
assertion that a change occurred. Section~\ref{sec:ai-guardrail-audit} of
Chapter~\ref{ch:ai-development} discusses two refusal codes specific to
one family of deliberately unfinished theorems ---
\texttt{COUNTERMODEL\_PROMOTION\_REFUSED} and its associated guard name,
\texttt{countermodel\_not\_promoted} --- as an instance of this same
family, added in direct response to an audited failure rather than
speculatively.

\section{Admission with a Covers Proof}
\label{sec:mfact-admission}

The move from candidate to artifact is not a boolean flag but a bundled
record. An \texttt{Admission} gathers the evidence for one candidate:
the kernel acceptance record, the sorry census, the axiom-audit output,
and a \texttt{covers} field --- a proof term witnessing that the
collected evidence entries jointly cover the obligations the admission
gate imposes. This last field is the structural device that prevents an
admission from being assembled with a silently missing check: the type
of \texttt{covers} forces every required obligation to be discharged by
some evidence entry before an \texttt{Admission} value can be
constructed at all, so ``we forgot to check the axiom footprint'' is not
a runtime bug this framework can express, only a compile-time
impossibility. The gate function, \texttt{admit}, is total: applied to a
candidate it returns either an \texttt{Admission} or a \texttt{Refusal},
with no third outcome such as a pending or ambiguous state. Totality
here is doing real epistemic work --- it is what makes ``the candidate
was neither admitted nor refused'' an unrepresentable claim rather than
merely a discouraged one.

\section{Manifest, Certified Release, and the Closed Objection Family}
\label{sec:mfact-manifest}

A \texttt{Manifest}, serializable to JSON, lists artifacts together with
their content hashes, axiom sets, fixture results, and the proven or
stated status of each declaration --- the object rendered, for this
release, into Table~\ref{tab:mfact-final-status} above and into the
per-declaration ladder discussed further in Chapter~\ref{ch:empirical}.
A \texttt{CertifiedRelease} pairs a manifest with its audit obligations.
Against a release \texttt{R}, \texttt{ValidObjection R} is a closed
inductive family whose constructors are the \emph{only} admissible forms
an objection to \texttt{R} may take --- an unaudited theorem, a missing
content hash, an unauthorized axiom, a failing fixture --- and each
constructor demands its own evidence to be instantiated, so an objection
cannot itself be raised without being checkable. The per-release theorem
\begin{center}
\texttt{theorem no\_valid\_objection : IsEmpty (ValidObjection R)}
\end{center}
is discharged by refuting each constructor of \texttt{ValidObjection}
against the corresponding checked field of \texttt{R}: there is no
unaudited theorem because the audit ran, no missing hash because the
manifest is complete, and so on for each constructor in turn. For the
release this thesis reports, the theorem is proven and, per
\texttt{\#print axioms}, depends on no axioms at all --- not even the
three standard Lean/Mathlib axioms (\texttt{propext},
\texttt{Classical.choice}, \texttt{Quot.sound}) that most of
\textsf{procint}'s own theorems carry. This is deliberately the
strongest possible case in the release: the framework's own
closed-objection theorem is proven from definitional unfolding alone.
The theorem's scope is exactly its constructors, and this thesis is
careful not to overstate it: \texttt{no\_valid\_objection} says the
\emph{enumerated} objection forms are refuted against this release's
checked fields, not that the release is beyond all possible criticism,
including criticism outside the closed family's four constructors (for
example, a citation-fidelity objection of the kind STANDING.md
explicitly flags as outside this pipeline's scope in its own Scope
section --- a limitation returned to in Chapter~\ref{ch:discussion}).

\subsection{The Boundary of Present Standing}
\label{sec:mfact-boundary}

Because status is algebra and not prose, this framework is also
positioned to state precisely, rather than gesture at, what remains
outside \textsc{Proven} in the current catalog. Of the current release's
\StatedCount{} stated declarations, two categories are worth naming
explicitly because they have opposite epistemic character. Five
declarations in the workflow countermodel family ---
\texttt{crownCounter\_reaches\_mid}, \texttt{crownCounter\_reaches\_final},
\texttt{crownCounter\_sound}, \texttt{crownCounter\_not\_bounded}, and
\texttt{WfNet.infinite\_transition\_countermodel\_sound\_not\_bounded} ---
are \emph{deliberately and permanently} stated: they witness why the
crown-jewel soundness equivalence (discussed in
Chapter~\ref{ch:procint}) requires the hypothesis \texttt{[Finite T]},
by exhibiting a workflow net with infinitely many transitions that is
sound but whose short-circuited net is unbounded. These are not
unfinished proof obligations queued for future work; they are a negative
witness, and Chapter~\ref{ch:ai-development} describes the mechanical
guard that refuses any attempt to promote them to proven standing. By
contrast, exactly one declaration in the current \TotalDecls-declaration
catalog is genuinely open proof work:
\texttt{ProcInt.BranchingProcess.isUnfoldingOf\_statement}, an
unfolding-correctness statement tracked separately in
\texttt{research/wfnet/obligations.toml} as obligation order~3
(\texttt{unfolding\_correctness}). Its statement, as recorded in
\texttt{packs/lean-math-pack/fragments/petriext.ttl}, attributes the
homomorphism condition it formalizes to Engelfriet's 1991 treatment of
branching-process unfoldings\footnote{The declaration's docstring reads,
verbatim: ``a branching process is a legal unfolding of net N when its
occurrence net is acyclic, has no backward conflict, and the labeling
respects the pre/post structure of N on singleton arcs (Engelfriet 1991,
Def.~3.1 homomorphism condition, specialized to ordinary nets).'' No
bibliographic entry for this attribution exists in \texttt{paper/refs.bib};
this thesis reports the attribution as the source states it, in a
footnote rather than as a \texttt{\textbackslash cite}, rather than
inventing a bibtex record to close the gap.}. This distinction --- a
deliberately closed negative witness versus a single, honestly reported
open item --- is itself an instance of the status algebra this section
opened with: neither is described here as a gap in the ordinary sense,
because only one of the two actually is one.

\section{The Ledger Law: Authority from the Manifest, Not the Path}
\label{sec:mfact-ledger}

A recurring temptation in a manufactured-artifact repository is to infer
a file's authority from where it lives or how it is named --- a
\texttt{generated/} directory, a header comment reading ``DO NOT
EDIT.'' This repository's ledger law rejects that inference explicitly.
There are, in its own words, no generated files; there are only
artifacts with receipts, and every repository file is first-class at its
canonical path. Authority comes from the artifact ledger,
\texttt{.mfact/artifacts.toml}, and from nowhere else. Concretely, this
has two directions of consequence. If a file is ledgered as produced by
ggen or a builder script, it may not be patched directly as a final
solution to a problem it exhibits: the sources or templates that produce
it must be modified, the file re-rendered, and \texttt{just regen-check}
must pass again before the edit is considered done --- an unreplayable
hand-edit to a ledgered file is \texttt{ARTIFACT\_DRIFT\_REFUSED}
regardless of how small or how obviously correct the hand-edit looked.
If, conversely, a file is \emph{not} ledgered but nonetheless carries
release standing, counts, audit status, or certification data, it is
classified as \texttt{ORPHAN\_ARTIFACT\_REFUSED}: the task is to either
ledger it properly or refuse to proceed, never to let an unregistered
file quietly carry standing.

Two surfaces are explicitly and permanently excluded from this ledger,
by design rather than by omission: \texttt{procint/Playground/**}, a
hand-authored Lean demonstration surface that is never ggen-rendered and
carries no standing claims of its own, and \texttt{pylab/**}, the
analogous hand-authored Python research surface. Both are developed at
length in Chapter~\ref{ch:ai-development}, which also describes a
mechanized scanner, itself living in \texttt{pylab/}, whose job is to
detect exactly the two ledger-law violations named in this section ---
drifted ledgered artifacts and unledgered orphans --- as a read-only,
repeatable check rather than a one-time manual audit.

\section{The Agent Actuation Constitution}
\label{sec:mfact-actuation}

Because candidates in this repository include agent-produced patches at
every layer --- TTL fragments, ggen templates, builder scripts, and this
thesis's own prose --- the repository constrains not just what may be
edited but how any change is permitted to take effect. Agents actuate
only through \texttt{just} recipes; raw invocations of \texttt{lake},
\texttt{ggen}, \texttt{mfact}, or git as a final actuation path are
disallowed unless a recipe or a doctor report explicitly instructs
otherwise for debugging. If a task requires an actuation path that does
not yet exist as a recipe, the constitution requires adding the recipe
first and only then using it, rather than reaching for the underlying
tool directly. This is not merely a style preference: a raw command
invoked outside its recipe bypasses whatever bookkeeping the recipe
performs around it (writing an ephemeral report, checking a
precondition, chaining a dependent step), reintroducing exactly the kind
of untracked side effect the ledger law in
Section~\ref{sec:mfact-ledger} exists to prevent.

The recipe surface itself is split into two classes with different
mutation rights. Diagnostic recipes --- \texttt{status}, \texttt{next},
\texttt{doctor}, \texttt{trace}, \texttt{why}, \texttt{theorem-status},
\texttt{proof-blockers}, \texttt{fixtures}, \texttt{docs-check} --- are
read-only: none of them may dirty the working tree, and none of their
output is itself a certified artifact. Only \texttt{report-write},
\texttt{check}, and \texttt{release} write the ephemeral cockpit reports
under \texttt{.mfact/reports/latest.*}, and these reports are themselves
gitignored and never ledgered --- they are a working surface for the
agent driving the recipe, not a standing claim. Certified status, by
contrast, lives only in the \texttt{release/} artifacts produced by the
manufacturing recipes proper: \texttt{render}, \texttt{build},
\texttt{audit}, \texttt{manifest}, \texttt{certify}, and
\texttt{standing-quadrature}. A completed task's report is required to
name the specific recipes it ran, not an ad hoc shell history, so that
the report itself is reconstructible from the recipe names alone.

\section{Core Release Identity and Recursion Refusal}
\label{sec:mfact-core-identity}

The manufacturing law of Section~\ref{sec:mfact-receipted-law} applies
recursively: the paper and this thesis are themselves artifacts of the
pipeline they describe, and their own claims must be receipted rather
than asserted. This creates a specific hazard the framework must guard
against explicitly. The core v26.7.7 release is frozen at tag
\texttt{\ReleaseTag}, with a manifest fold hash, a proven-declaration
count, and a total-declaration count fixed at that tagged commit ---
values this thesis reports throughout via the \texttt{release\_macros}
fragment rather than by typing them by hand. Post-release work,
including work on this thesis and on the correspondence and evaluation
material of later chapters, must not mutate the identity of the core
certified release it reports on. If an operation would change the core
fold hash, the proven count, the declaration count, or the manifest
itself, the constitution requires one of two outcomes and forbids the
third: either refuse the operation as
\texttt{RECEIPT\_RECURSION\_REFUSED}, or promote it explicitly into a
new core certification cycle with new manifest values recorded under a
new identity. What is not permitted is a silent third option in which
the core identity shifts underneath an unchanged tag. Post-release
witnesses such as \texttt{ProcInt.Release.PostRelease} are counted
separately, as \texttt{POST\_RELEASE\_WITNESSES}, and are never folded
into the core proven count; the core tag itself is pinned to the core
release commit and is not moved by later packet work, which receives its
own tags and its own packet hashes. \texttt{just certify} and
\texttt{just release} are required to verify that the release commit is
an ancestor of, or equal to, the certified tag before reporting a
passing \texttt{CERTIFIED\_RELEASE} status; a tag that predates the
currently rendered artifacts is treated as
\texttt{RECEIPT\_RECURSION\_REFUSED}, not filed away as a historical
curiosity to be noted and passed over. Chapter~\ref{ch:ai-development}
describes the mechanized check that implements this ancestry
verification and traces it to the specific 2026-07-07 audit finding that
motivated adding it.

\section{Summary}

This chapter has developed \textsf{mfact} as the typed apparatus that
makes the pipeline's receipted manufacturing law operational: candidates
and refusals as closed data types with no silent failure path;
admission as a total function whose \texttt{covers} proof makes an
incomplete check unrepresentable; the manifest, certified release, and
closed \texttt{ValidObjection} family as the objects a
\texttt{no\_valid\_objection} theorem is proven against, scoped exactly
to its enumerated constructors; the ledger law that locates authority in
\texttt{.mfact/artifacts.toml} rather than in file paths, headers, or
naming conventions; the agent actuation constitution that confines any
producer's actuation to named, classified \texttt{just} recipes; and the
core release identity's protection against self-referential mutation.
Chapter~\ref{ch:procint} now turns to the first substantial artifact
this apparatus was used to manufacture: \textsf{procint}, and in
particular the workflow-net soundness development whose crown-jewel
theorem and deliberate countermodel were introduced only by name in
Section~\ref{sec:mfact-boundary} above.

\chapter{procint: Process Intelligence as Process Law}
\label{ch:procint}
\label{sec:procint}

Chapter~\ref{ch:mfact} developed \textsf{mfact} as a content-neutral
admission apparatus: candidates, typed refusals, admissions with covers
proofs, manifests, and the closed family of valid objections say nothing,
by themselves, about what is being admitted. This chapter turns to the
first domain manufactured through that apparatus: \textsf{procint}, a
Lean~4 formalization of process-intelligence semantics --- Petri nets,
workflow nets and their soundness property, alignment-based conformance
checking, and OCEL~2.0 object-centric event logs. \textsf{procint} is not
a demonstration corpus assembled to exercise \textsf{mfact}; it is the
substrate the rest of this thesis's empirical claims are about, and its
declaration catalog, drafted across independent LLM-assisted lanes and
admitted only through kernel testing, is exactly the kind of
manufactured-not-authored content the mfact gate exists to police. As
the accompanying paper's ``procint: Process Intelligence as Process Law''
section puts it, process
intelligence in this sense does not mean dashboarding over event logs; it
means a mathematical substrate for deciding, replaying, and explaining
process standing --- what happened, what was admitted, what was refused,
which object relations were involved, which conformance obligations
passed, and which receipts authorize the resulting claim.

\textsf{procint} depends on Mathlib~\cite{mathlib2020} for its background
mathematics and on CSLib~\cite{cslib2026} for transition-system and trace
infrastructure, and is organized into namespaces that separate concerns
cleanly along the process-intelligence canon: \texttt{ProcInt.Foundations}
for metrics, identifiers, and labeled transition systems;
\texttt{ProcInt.Petri} for nets, firing, reachability, the state
equation, invariants, boundedness, liveness, object-centric Petri nets,
stochastic nets, and unfoldings; \texttt{ProcInt.Workflow} for workflow
nets, the short-circuit construction, and soundness; \texttt{ProcInt.Logs}
for events, traces with monotone timestamps, event logs, and XES;
\texttt{ProcInt.Models} for process trees, POWL, directly-follows graphs,
causal nets, BPMN, and Declare; \texttt{ProcInt.Conformance} for moves,
alignments, cost, fitness, and token replay; \texttt{ProcInt.Ocel} for
OCEL~2.0 events, objects, event-to-object and object-to-object relations,
flattening, and lifecycles; \texttt{ProcInt.Ocpq} for object-centric
process queries; \texttt{ProcInt.Analytics} for temporal, causal,
correlational, cube, and streaming analytics; and \texttt{ProcInt.Registry}
for the cross-product algorithm and model-breed registries that are
themselves data-driven rows of the same declaration graph rather than
hand-written tables. This chapter takes the workflow-net soundness
development as its center of gravity, because it is where the corpus's
single crown-jewel theorem lives, but treats the surrounding layers ---
Petri semantics beneath it, conformance and OCEL~2.0 beside it --- with
enough precision to support the canon-coverage claims the accompanying
paper's canon-status table makes.

\section{Petri Net Semantics}
\label{sec:procint-petri}

The base layer of \textsf{procint}'s process semantics follows Murata's
survey treatment of Petri nets~\cite{murata1989}, formalized directly
rather than via an intermediate encoding. A net is places, transitions,
and weighted flow over finite types, and a marking is a finitely
supported function \(M : P \to_{f} \mathbb{N}\) rather than a bare
function into the naturals; representing markings as \texttt{Finsupp}
values, Mathlib's finitely supported function type, does two things at
once. First, it unifies the weighted-net and colored-net cases: a
weighted marking is a \texttt{Finsupp} into \(\mathbb{N}\), and richer
codomains extend the same representation without a second definitional
layer. Second, it gives the whole marking algebra --- firing, the state
equation \(M' = M + A^{\mathsf T}u\) relating a reached marking to the
initial marking via the incidence matrix \(A\) and a firing-count vector
\(u\), and place/transition invariants --- as lemmas already available
over \texttt{Finsupp} in Mathlib, with truncated subtraction standing in
for the natural-number semantics of ``consume what is present.'' This
representational choice recurs later in this chapter: the same
truncated-subtraction discipline is what the crown-jewel theorem's
countermodel exploits, and what the correspondence rail of
Chapter~\ref{ch:correspondence} has to reconcile against Rust's own
saturating arithmetic.

Firing, reachability, boundedness, and liveness are defined over this
marking algebra in \texttt{ProcInt.Petri} and reused, rather than
redefined, by the workflow-net layer above it and by the object-centric
Petri net (OCPN) and stochastic-net extensions that follow
the accompanying paper's canon-status table's ``formalized''/``seed
instance'' distinction: the
core marking, firing, and invariant theory is formalized in full, while
OCPN and stochastic nets carry seed instances built on top of it rather
than independent full developments.

\section{Workflow Nets and the Soundness Property}
\label{sec:procint-wfnet}

A \texttt{WfNet} carries a Petri net as a field --- composition, not
inheritance by extension --- together with the source-place, sink-place,
and connectivity conditions of van der Aalst's Definition~6--7
\cite{vanderaalst1997}: a single source place with no input transitions, a
single sink place with no output transitions, and every node on some path
from source to sink. This composition-over-extension choice keeps the
underlying Petri net's marking algebra directly reusable rather than
re-derived under a workflow-specific wrapper.

The classical soundness property of a workflow net is usually stated as
one conjunction: from the initial marking, the final marking is always
reachable (option to complete); when the final marking is reached, no
tokens remain elsewhere (proper completion); and every transition can
still fire from some reachable marking (no dead transitions). Following
the observation that these three clauses are logically independent of one
another, \textsf{procint} does not encode soundness as a single
conjunction-shaped definition. Instead it states option-to-complete,
proper-completion, and no-dead-transitions as three separately named
predicates, with \texttt{Sound} defined as their explicit conjunction.
This keeps partial results honest in a literal, checkable sense: a lemma
that establishes one clause is stated, and can be cited, about exactly
that clause, and the fact that \S\ref{sec:procint-crown} below can report
\texttt{PROVEN\_SUPPORT} status for two of the three clauses separately
from the equivalence that uses all three depends directly on this
decomposition.

\section{The Crown-Jewel Theorem}
\label{sec:procint-crown}

The single most load-bearing result in the \textsf{procint} corpus is van
der Aalst's soundness characterization, Lemma~8 of
\cite{vanderaalst1997}: for a workflow net satisfying the source/sink/
connectivity conditions above, \texttt{Sound} holds of the net if and only
if the short-circuited net --- the net obtained by adding a transition
from the sink place back to the source place, folding the ``process
completes and a new instance may start'' behavior into an ordinary
liveness-and-boundedness question --- is both live (every transition can
recur from every reachable marking) and bounded (the reachable marking
set is finite). \textsf{procint} calls this the crown-jewel theorem
because the release's Standing Quadrature and manifest apparatus
(Chapter~\ref{ch:empirical}) is organized around it as the single result
whose kernel-admitted status the release most directly advertises, and
because it is the result that most clearly demonstrates the pipeline
producing a theorem whose informal target predates the pipeline by three
decades, rather than a theorem invented to be easy to formalize.

\begin{theorem}[Crown-jewel soundness equivalence, informal statement]
\label{thm:crown-informal}
Let \(W\) be a workflow net over a finite transition type \(T\) satisfying
the source, sink, and on-path connectivity conditions of van der Aalst's
Definition~6--7. Then \(W\) is sound if and only if the short-circuited
net \(W^{sc}\) is live and bounded.
\end{theorem}

The obligation ladder recorded in the repository's crown research
tracking (\texttt{research/wfnet/obligations.toml}) states the run's
standing toward this theorem as four ordered obligations. Order~0 is the
equivalence itself, annotated \emph{``the crown equivalence itself (van
der Aalst 1997, Lemma~8)''} and carrying status \texttt{PROVEN}. Orders~1
and~2, \emph{``soundness implies the final marking is reachable''} and
\emph{``soundness implies every transition can fire,''} carry status
\texttt{PROVEN\_SUPPORT}: they are the proper-completion and
no-dead-transitions directions of the soundness decomposition from
\S\ref{sec:procint-wfnet}, each proven separately and feeding the
equivalence rather than duplicating it. Order~3, discussed on its own in
\S\ref{sec:procint-open} below, is the one obligation in this ladder ---
and, as \S\ref{sec:procint-open} details, in the entire declaration
catalog --- that remains \texttt{STATED} rather than \texttt{PROVEN}. The
\texttt{[crown]} block of the same file records the release-level
judgment, \texttt{status = "PROVEN"}, computed the same way every other
status claim in this thesis is computed: by grepping the admitted TTL
declaration catalog for the relevant declarations' recorded status, never
by hand-asserting the release's headline claim. The generated fragment
\texttt{crown\_jewel\_status.tex} (quoted in full in
Chapter~\ref{ch:ai-development})
reports this same ground-truth classification directly from
\texttt{ProcInt.crownJewel\_status}, so this chapter's prose claim and the
release's generated claim are, structurally, the same read of the same
underlying value rather than two independently maintained assertions that
could drift apart.

\subsection{The \texttt{[Finite T]} Hypothesis and Its Countermodel}
\label{sec:procint-finite-t}

Theorem~\ref{thm:crown-informal} as stated above is not unconditionally
true of van der Aalst's original formulation once the transition set is
allowed to be infinite, and \textsf{procint}'s Lean statement of the
crown-jewel theorem is correspondingly narrower than the informal
statement: it carries an explicit \texttt{[Finite T]} instance hypothesis
on the transition type. This hypothesis is not a formalization
convenience adopted for proof-engineering ease; it is necessary, and the
repository carries a Lean-admitted witness proving that it is necessary
rather than merely asserting so in prose.

The witness is a purpose-built counterexample net, assembled in five
stages across the \texttt{Workflow.Countermodel} module (declared in
\texttt{packs/lean-math-pack/fragments/workflow\_countermodel.ttl}, module
order 34) and structured around a place type
\texttt{CrownCounterPlace} (an input place \texttt{i}, a queue place
\texttt{q}, an output place \texttt{o}, and a countably infinite family of
counter places \texttt{c n}) together with a transition type
\texttt{CrownCounterTransition := ℕ ⊕ ℕ}, one summand for an
\emph{absorb} transition at each index \(n\) and one for a matching
\emph{emit} transition. Absorb transition \(n\) consumes one token from
\texttt{i} and produces \(n\) tokens into \texttt{q} together with one
token into \texttt{c n}; emit transition \(n\) consumes those \(n\) tokens
from \texttt{q} and the one token from \texttt{c n}, and produces one
token into \texttt{o}. Because \(\mathbb{N} \oplus \mathbb{N}\) is
infinite --- witnessed directly by the instance declaration
\texttt{infinite\_CrownCounterTransition} --- this net has infinitely many
transitions, one absorb/emit pair per natural number, and can therefore
never satisfy \texttt{[Finite T]}.

The construction is engineered so that this infinite net is nonetheless
sound while its short-circuit is unbounded, the exact combination the
crown-jewel equivalence rules out under \texttt{[Finite T]}. Five
declarations carry this argument: \texttt{crownCounter\_reaches\_mid} and
\texttt{crownCounter\_reaches\_final} establish the reachability facts the
soundness argument needs; \texttt{crownCounter\_sound} establishes that
the net satisfies all three soundness clauses of
\S\ref{sec:procint-wfnet}; \texttt{crownCounter\_not\_bounded} establishes
that for every bound \(k\), some reachable marking of the short-circuit
violates it, because absorbing at index \(n\) drives \(n\) tokens into
\texttt{q} for arbitrarily large \(n\); and the main theorem,
\texttt{WfNet.infinite\_transition\_countermodel\_sound\_not\_bounded},
assembles these into the conjunction \(\mathrm{Infinite}\,
\text{CrownCounterTransition} \land \exists\, W,\, W.\mathrm{Sound}
\land \lnot \exists\, k,\, W.\mathrm{shortCircuit}.\mathrm{Bounded}\,
W.\mathrm{initialMarking}\,k\).

These five declarations are recorded in the catalog with \texttt{status
"stated"}, and this is a deliberate, permanent classification rather than
an unfinished proof obligation awaiting a future lane. The countermodel's
job is to witness a negative fact --- that the equivalence fails outside
\texttt{[Finite T]} --- and the repository's guardrails treat promoting
these five declarations to a proven, axiom-audited status as a specific
failure mode to be mechanically blocked rather than merely discouraged:
a countermodel family that quietly acquired \texttt{PROVEN} status would
misrepresent what the family is for. The corresponding refusal vocabulary
--- \texttt{WFNET\_INFINITE\_TRANSITION\_COUNTERMODEL} as the status key
for this theorem family, \texttt{countermodel\_not\_promoted} as the guard
name, and \texttt{COUNTERMODEL\_PROMOTION\_REFUSED} as the refusal fired
if a promotion is attempted --- exists precisely so that this
countermodel's status cannot silently drift from ``deliberate negative
witness'' to ``forgotten stub.'' Read correctly, then, these five
declarations are not a gap in the crown-jewel development; they are part
of its statement, in the same sense that a theorem's hypotheses are part
of its statement. \texttt{[Finite T]} is not a hedge adopted for proof
convenience --- it is exactly the hypothesis the countermodel shows is
required, made visible in the formal statement rather than left implicit
in prose the way an informal restatement of Lemma~8 might leave it.

\section{Conformance Checking and Alignments}
\label{sec:procint-conformance}

\textsf{procint}'s conformance-checking layer, \texttt{ProcInt.Conformance},
follows Carmona et al.'s alignment-based treatment of conformance checking
\cite{carmona2018} in preference to naive token-replay diagnostics.
Alignments relate an observed trace to a model run via four move kinds:
synchronous moves where the log and the model agree, log-only moves,
model-only moves, and silent (\(\tau\)) moves; a cost function over moves
induces optimal alignments as the alignments of least total cost. Fitness
--- the conformance metric this thesis returns to in
Chapter~\ref{ch:correspondence}, where it is the quantity a Rust-extracted
implementation and a hand-authored Lean model are related over --- is
represented as a \texttt{Between01} value by construction, so that the
boundedness of the metric (fitness always lies in \([0,1]\)) is a
type-level fact enforced by the representation itself rather than a
theorem that has to be separately proven and could, in principle, have
been forgotten.

\section{OCEL 2.0 Coverage}
\label{sec:procint-ocel}

Object-centric event logs follow the OCEL~2.0 specification
\cite{ocel2}, Definitions~1--2, in \texttt{ProcInt.Ocel}: typed events and
typed objects, event-to-object and object-to-object relations carrying
qualifiers, and attribute maps over both events and objects. Flattening
onto a classical, single-case-notion event log, and object lifecycle
extraction, are represented as derived operations over this base
representation rather than as independently axiomatized structures, so
that a flattening result is a theorem about the OCEL~2.0 representation
rather than an assumption stated separately about a flattened log.
Object-centric Petri nets, the net-level counterpart of object-centric
event data, follow~\cite{ocpn2020} and are formalized, per
the accompanying paper's canon-status table, as a seed instance built
on the Petri semantics of \S\ref{sec:procint-petri} rather than as a full
independent development at the same depth as the workflow-net soundness
theory.

\section{The One Open Obligation: Unfolding Correctness}
\label{sec:procint-open}

Of the \TotalDecls\ declarations recorded in the catalog at the time of
the release described by \texttt{release\_macros.tex},
exactly one is genuinely open proof
work rather than a completed theorem, a definition, a test oracle, or a
deliberately stated negative witness of the kind discussed in
\S\ref{sec:procint-finite-t}: order-3 of the crown obligation ladder,
\texttt{ProcInt.BranchingProcess.isUnfoldingOf\_statement}, tracked in
\texttt{research/wfnet/obligations.toml} under the name
\texttt{unfolding\_correctness} with status \texttt{STATED} and annotated
\emph{``branching-process unfolding statement (separate stated lane).''}
This is worth stating plainly rather than folding into the crown-jewel
narrative above, because it is a different kind of incompleteness from
the countermodel family: the countermodel is deliberately, permanently
stated as a negative witness that should never be promoted; the unfolding
statement is stated because the corresponding proof has not yet been
produced, and promoting it to \texttt{proven} remains an open task for a
future lane, not a status this repository's guardrails forbid in
principle.

The declaration, recorded in module \texttt{Petri.Unfolding} in
\texttt{packs/lean-math-pack/fragments/petriext.ttl}, states --- but does
not yet prove --- what it means for a branching process to be a legal
unfolding of a Petri net \(N\): given a branching process \(B\) over
places \(P\), transitions \(T\), condition labels \(C\), and events \(E\),
\(B\) is an unfolding of \(N\) exactly when \(B\)'s occurrence net is
acyclic, has no backward conflict, and its event/condition labeling
respects \(N\)'s pre- and post-set structure on singleton arcs:
\begin{center}
\begin{tabular}{l}
\(B.\mathrm{occ.Acyclic} \;\land\; B.\mathrm{occ.NoBackwardConflict}\)\\
\(\land\; (\forall\, e\, c,\; c \in B.\mathrm{occ.preC}\,e \to N.\mathrm{pre}\,(B.\mathrm{eventTrans}\,e)\,(B.\mathrm{condPlace}\,c) \neq 0)\)\\
\(\land\; (\forall\, e\, c,\; c \in B.\mathrm{occ.postC}\,e \to N.\mathrm{post}\,(B.\mathrm{eventTrans}\,e)\,(B.\mathrm{condPlace}\,c) \neq 0)\)
\end{tabular}
\end{center}
The declaration's own doc comment attributes the homomorphism condition
this statement specializes --- to ordinary nets, on singleton arcs --- to
Engelfriet's Definition~3.1 characterization of branching-process
unfoldings (Engelfriet 1991), as the source fragment states it directly.
No bibliographic entry for this attribution exists in this thesis's
reference list, and none is added here: inventing a bibtex record for a
citation that appears only as an in-source doc-comment attribution, absent
independent verification of its bibliographic details, would be exactly
the kind of fabrication this thesis's own manufacturing discipline is
designed to refuse elsewhere in the pipeline, and prose is not exempt
from that discipline merely because it is prose rather than a TTL
declaration. The attribution is therefore reported here as the source
states it --- an in-repository doc comment, not a peer-reviewed citation
this thesis independently verifies --- and readers who need the primary
source should consult the branching-process unfolding literature
directly rather than treat the name and year given here as a full
bibliographic record.

Because \texttt{isUnfoldingOf\_statement} is a \texttt{def} of type
\texttt{Prop} rather than a \texttt{theorem} with a completed proof term,
it participates in the corpus the way any stated-but-undischarged
statement does: it is visible in the declaration catalog, it is counted
among the \StatedCount\ stated declarations reported in
Table~\ref{tab:corpus} of \texttt{release\_macros.tex}, and it is excluded from the
\KernelTheorems\ kernel-admitted, axiom-audited theorems the release
advertises as proven. Its presence in the catalog at \texttt{STATED}
status, rather than its silent omission, is itself part of what this
thesis means by an honest, receipted release: a formalization effort that
reports only its completed theorems and omits its open statements would
overstate its own standing, and the manifest-and-quadrature apparatus of
Chapter~\ref{ch:empirical} is built specifically to make that kind of
omission structurally visible rather than merely discouraged in review.

\chapter{The Correspondence Rail: Bridging Rust and Lean}
\label{ch:correspondence}

Chapter~\ref{ch:procint} manufactured domain semantics --- Petri nets,
workflow-net soundness, conformance, OCEL~2.0 --- from a declared
obligation graph whose objects are entirely Lean statements about Lean
structures. This chapter develops a second, structurally different
manufacturing rail opened by \textsf{mfact}: implementation
correspondence, in which the obligation graph relates a Lean declaration
not to another Lean declaration but to code extracted from an external,
unverified Rust crate, \texttt{wasm4pm-core}. Where the procint
formalization's rail (Chapter~\ref{ch:procint}) manufactures theorems
purely within Lean's own world of hand-written
types, this rail manufactures theorems that reach outside that world to
quantify over a generated Lean image of Rust source the pipeline itself
did not author and does not trust by default. This chapter's central
worked example, the D1 token-replay-counts correspondence obligation, is
also this thesis's clearest illustration of one of the guardrails
introduced after the post-\texttt{v26.7.7} five-rail audit: a
correspondence theorem that fails to quantify over its actual extracted
declaration is not a weaker version of the intended result, it is a
different, insufficient result, and the repository's own history records
a case where exactly this distinction was caught and corrected before the
obligation was allowed to report proven status.

\section{Aeneas, Charon, and Functional Extraction}
\label{sec:correspondence-aeneas}

Aeneas~\cite{ho2022aeneas} is a tool for verifying Rust programs by
translating them into a purely functional representation suitable for a
proof assistant, rather than by reasoning about Rust's operational
semantics directly; its Charon front end~\cite{protzenko2024charon}
separates this translation into two stages, first recovering a clean
syntactic and typed representation of the Rust source (Charon's job),
then performing the semantic translation into the functional target
proper (Aeneas's job). Aeneas is designed to target whole Rust programs
under a full memory model, including borrowing and mutation; the use of
it in this repository is narrower by design, applied to a single,
dependency-free arithmetic module of \texttt{wasm4pm-core} rather than to
a full program, because --- consistent with the untrusted-producer
discipline this thesis applies throughout --- a single bounded instance
of the rail, fully receipted, is what standing requires before the rail
is generalized to a wider Rust surface, not a demonstration that Aeneas
can be pointed at arbitrary Rust code.

Extraction's output is, by the same discipline that governs every other
producer in this pipeline, a candidate, never a proof. A Charon/Aeneas
run over \texttt{wasm4pm-core/src/conformance\_counts.rs} yields a Lean
module of extracted, proof-free declarations --- the
\texttt{ReplayCounts\{produced, consumed, missing, remaining\}} struct and
the integer \texttt{fitness\_num}/\texttt{fitness\_den} functions, taken
verbatim from the crate's arithmetic --- and this extracted image
acquires no standing by virtue of having been extracted. Standing still
requires kernel admission of a hand-authored correspondence theorem
stated \emph{over} the extracted declaration: an abstraction function
(\texttt{Abs.toSpec} in the repository's own naming) lifts the extracted,
proof-free struct into \textsf{procint}'s own \texttt{ReplayCounts} type
--- the same type introduced in \S\ref{sec:procint-conformance}, which
carries \texttt{missing} \(\le\) \texttt{consumed} and \texttt{remaining}
\(\le\) \texttt{produced} as proof fields, not merely as informal
invariants --- given exactly those two inequalities as hypotheses. This
is the general shape of every obligation on this rail: extraction
produces data, the abstraction function is where the correspondence claim
is actually made, and the correspondence theorem is where that claim is
kernel-checked. Section~\ref{sec:correspondence-guard} returns to why the
Lean statement's literal dependence on the extracted declaration, rather
than on an independently hand-written stand-in of the same shape, is the
part of this rail a guard specifically checks for.

\section{The D1 Obligation Family and \texttt{ValidObjection}}
\label{sec:correspondence-d1}

The correspondence rail's status ladder --- declared, extracted, stated,
proven --- mirrors the proven/stated distinction Chapter~\ref{ch:procint}
uses throughout, but is computed independently, by a dedicated builder
(\texttt{scripts/build\_verif.py}) reading observed evidence rather than
by hand: a Charon/Aeneas pipeline receipt establishing that extraction
actually ran and produced the declaration a correspondence theorem
claims to quantify over, a \texttt{lake build} of the extracted module
against \textsf{procint} establishing that the extracted image and the
hand-authored model actually type-check together, and an axiom-print
check on the resulting theorem establishing that its proof, if one
exists, does not smuggle in an unaudited axiom. The rail's first
obligation, and to date its only instance, is scoped deliberately
narrowly: token-replay counts correspondence, obligation ``D1'' in the
repository's own numbering, relating the extracted
\texttt{fitness\_num}/\texttt{fitness\_den} integer functions to
\textsf{procint}'s \texttt{ReplayCounts}/\texttt{fitness} model
introduced in \S\ref{sec:procint-conformance}. The obligation does not
assert that the crate's integer ratio and \textsf{procint}'s exact-rational
fitness are numerically equal in general --- they are structurally
different formulas, a single integer fraction against an average of two
ratios --- only that each computes correctly from the same underlying
counts; forcing a false numeric equality between them was rejected in
favor of this honest, weaker, but kernel-provable statement, a choice
\S\ref{sec:correspondence-d1-history} below traces through the
repository's own record of arriving at it.

This obligation family is exactly the kind of claim
\texttt{ValidObjection R}, introduced in Chapter~\ref{ch:mfact} as the
closed inductive family of admissible objection forms against a
certified release \(R\), is positioned to police once the correspondence
rail's evidence is folded into a release manifest: an unaudited
correspondence theorem, a missing extraction receipt, or an unauthorized
axiom in the correspondence proof are each instances of the same
objection shapes \texttt{ValidObjection} already enumerates for the
\textsf{procint} corpus, not a separate objection vocabulary invented for
this rail. The correspondence rail does not relax the admission
discipline of Chapter~\ref{ch:mfact}; it applies the same discipline to a
second class of obligation, one whose statements happen to quantify over
extracted rather than hand-written declarations.

\section{A Worked Case: the Honest D1 Restatement}
\label{sec:correspondence-d1-history}

The repository's own record of this obligation's history, preserved in
\texttt{docs/HONEST\_D1\_STATEMENT.md}, is a documented instance of the
guardrail this chapter is centrally about, and is worth reconstructing in
some detail because it shows the failure mode being caught in the
repository's own past, not merely described in the abstract.

An early, placeholder version of the D1 theorem, appropriate to
\texttt{DECLARED} status but not to any status beyond it, took the shape
\begin{center}
\small
\begin{tabular}{l}
\texttt{theorem token\_replay\_counts\_corr}\\
\texttt{\quad(c : ProcInt.ReplayCounts) (num den : ℕ) (hden : den ≠ 0)}\\
\texttt{\quad(hnum : (num : ℚ) / (den : ℚ) = ProcInt.fitness c) :}\\
\texttt{\quad ProcInt.fitness c = (num : ℚ) / (den : ℚ) := by sorry}
\end{tabular}
\end{center}
This statement is tautological given its own hypothesis: it \emph{assumes}
the very equality between the extracted ratio and \textsf{procint}'s
fitness that the obligation is supposed to establish, then proves the
symmetric restatement of that assumption. As a declared placeholder ---
marking that the obligation exists and has a target shape, before any
extraction evidence is available --- this was an honest artifact of where
the obligation stood. Shipped as a \texttt{STATED} or \texttt{PROVEN}
claim about the actual correspondence between Rust and Lean fitness, it
would not have been: proving a restatement of an assumed equality
establishes nothing about whether the assumption itself holds of the
extracted code.

The repository's own record documents the honest correction made instead
of promoting this placeholder. \texttt{docs/HONEST\_D1\_STATEMENT.md}
records that, as of 2026-07-07, the obligation was rewritten to remove
the false assumption and replace it with two source hypotheses tracing
the extracted integer values back to \texttt{ReplayCounts} fields
(\texttt{num = c.consumed - c.missing}, \texttt{den = c.produced +
c.remaining}), with a conclusion stating both fitness formulas exactly ---
\textsf{procint}'s Rozinat-and-van-der-Aalst-style average of two ratios,
and the crate's single-fraction ratio --- as a conjunction, rather than
asserting they are equal. At that point in the obligation's history, the
document records its \texttt{verif:aeneasDecl} field, which is meant to
name the actual extracted Lean declaration the theorem quantifies over,
as the literal placeholder string \texttt{"TBD"}: the honest theorem
statement had been written, but the binding to a real Charon/Aeneas
extraction output had not yet been completed, so the theorem, however
honestly stated, could not yet be said to be \emph{about} extracted code
in the sense the rail requires. This is precisely the situation the
guardrail set out in \texttt{AGENTS.md} names directly: ``a receipt field
like \texttt{aeneasDecl: "TBD"} is itself a refusal signal (the binding
was never completed) and must block PROVEN status, not ship alongside
it.'' The document's own status table records the obligation, at that
point, as \texttt{STATED} with its proof body a \texttt{sorry} --- an
honest report of an incomplete binding, not a claim dressed up as more
complete than it was.

\section{Completing the Binding}
\label{sec:correspondence-guard}

The obligation as currently recorded in
\texttt{packs/lean-math-pack/fragments/verif.ttl} no longer carries a
\texttt{"TBD"} \texttt{aeneasDecl} field: the field now names
\texttt{ReplayCounts}, the actual extracted structure, and the
theorem's own statement has changed to match, quantifying directly over
\texttt{Wasm4pmVerify.Generated.ReplayCounts} --- the Charon/Aeneas
extraction target named by \texttt{verif:aeneasModule} as
\texttt{Wasm4pmVerify.Generated} --- rather than over a bare
\texttt{num}/\texttt{den} pair related to it only by hypothesis. The
current theorem takes a value \texttt{gen} of the extracted
\texttt{Wasm4pmVerify.Generated.ReplayCounts} type together with the two
proof obligations \texttt{gen.missing.val ≤ gen.consumed.val} and
\texttt{gen.remaining.val ≤ gen.produced.val}, applies the abstraction
function \texttt{Wasm4pmVerify.toSpec} to obtain a \textsf{procint}
\texttt{ReplayCounts} value \texttt{spec}, and concludes a conjunction:
that the abstraction preserves the two proof-field witnesses
(\texttt{spec.missing\_le = hm}, \texttt{spec.remaining\_le = hr}), that
\textsf{procint}'s fitness at \texttt{spec} computes to the
Rozinat-and-van-der-Aalst average-of-ratios formula, that the extracted
\texttt{fitness\_num} computes via saturating subtraction over the
extracted fields, and that the extracted \texttt{fitness\_den} either is
undefined or computes to some explicit value. This is exactly the
completed form of the honest statement
\S\ref{sec:correspondence-d1-history} traces the repository toward: the
extraction binding named \texttt{"TBD"} is filled in, the statement
quantifies over the real extracted declaration rather than over
independently hand-written stand-ins of the same shape, and the guardrail
this section is built around --- that a correspondence obligation
reported as proven must actually import and quantify over its extraction
via a declared abstraction function, not merely resemble it --- is
satisfied by construction rather than by assertion. The live evidence for
this obligation's current status, computed by \texttt{scripts/build\_verif.py}
from the observed charon/aeneas/lake evidence trail rather than typed by
hand into this chapter, is
\input{../paper/correspondence_status}
This table is regenerated by the same builder each time the underlying
evidence changes; this chapter reports its structure and history, not its
current numeric content, which is authoritative only in the rendered
fragment above.

\chapter{Empirical Evidence: praxis-graphlaw and rslab}
\label{ch:empirical}

Chapter~\ref{ch:procint} and Chapter~\ref{ch:correspondence} developed
two manufacturing rails whose evidence is, in the end, a Lean~4 term
checked by the kernel: a declaration either type-checks against its
stated proposition under the authorized axiom set, or it does not. This
chapter turns to a rail of a structurally different kind, one that
produces no kernel term at all. praxis-graphlaw is a Rust crate in the
sibling \texttt{praxis} workspace, and rslab is the evidence pipeline
that turns its raw command output into receipted, citable measurement.
Neither rail is a proof engine, and neither is presented in this thesis,
or in \texttt{paper/main.tex}, as one. The purpose of this chapter is to
state precisely what kind of evidence praxis-graphlaw and rslab do
produce, how that evidence is receipted, and why it is kept apart from
\textsf{procint}'s kernel-checked standing rather than blended into it.

\section{praxis-graphlaw: Executable Law-State Evaluation}
\label{sec:emp-praxis-graphlaw}

praxis-graphlaw is the executable law-state evaluation engine living in
the sibling \texttt{praxis} Rust workspace. It evaluates SHACL, ShEx, N3
forward-chaining, and Datalog dialects against RDF graphs, and exposes an
admission context, \texttt{bcinr\_powl::admit}, that gates graph mutation
on the outcome of those checks: a mutation is admitted only if the
relevant law-state evaluation passes. \texttt{.agents/explorer\_step3/handoff.md}
records the workspace \texttt{praxis} sits in as a Cargo workspace of
eleven further crates alongside \texttt{praxis-graphlaw} itself ---
\texttt{agent8}, \texttt{ggen}, \texttt{chatman-common},
\texttt{praxis-core}, \texttt{praxis-proposer}, \texttt{praxis-retrofit},
\texttt{rust-fable-testbed}, \texttt{powl2-decompose},
\texttt{pddl-index}, and \texttt{praxis-synthesis},
\texttt{praxis-lean} --- of which \texttt{ggen} is the same rendering
engine this thesis's \textsf{mfact} rail invokes to project the
\textsf{procint} catalog into Lean source (Chapter~\ref{ch:mfact}), and
\texttt{praxis-lean} names the crate in that workspace nearest to a Lean
integration surface, though this chapter's evidence concerns
\texttt{praxis-graphlaw} specifically, not that surface.

praxis-graphlaw is, in the vocabulary this thesis has used throughout, a
distinct manufacturing rail from \textsf{procint}. Where \textsf{procint}
formalizes process-mining semantics as Lean propositions and admits
theorems about them through the kernel (Chapter~\ref{ch:procint}),
praxis-graphlaw \emph{executes} law-state checks at runtime in Rust; it
is not admitted through any proof obligation, and no claim in this thesis
treats a passing SHACL, ShEx, N3, or Datalog evaluation as a
kernel-checked fact about the evaluator's own correctness. The two rails
are related only by citation: this chapter, following \texttt{paper/main.tex},
reports both, but neither certifies the other. In particular,
\textsf{mfact} does not claim, and this thesis does not assert, that
praxis-graphlaw's SHACL/ShEx/N3/Datalog evaluation logic is formally
verified in the sense Chapter~\ref{ch:mfact} and Chapter~\ref{ch:procint}
use that phrase for \textsf{procint} declarations. Its correctness
evidence is empirical: the conformance and end-to-end admission-gate test
suite receipted through the rslab rail described next, together with a
benchmark suite whose caveats are reported alongside the numbers rather
than smoothed over (Section~\ref{sec:emp-caveats}).

\emph{One further use of the word ``receipt'' recurs in this workspace and
is worth disambiguating up front.} \texttt{bcinr\_powl}'s admission
context governs a single graph-mutation decision at runtime, inside
praxis-graphlaw's own process; it is not the release-manifest receipt
chain of Chapter~\ref{ch:mfact}, and a passing admission decision inside
\texttt{bcinr\_powl} confers no \textsf{mfact} standing on anything.

\section{rslab: An Empirical Evidence Rail, Not a Proof Engine}
\label{sec:emp-rslab}

rslab's own doctrine, stated verbatim in \texttt{rslab/README.md} and
reproduced in \texttt{paper/main.tex}, separates three tiers of activity
that this thesis's manufacturing vocabulary keeps distinct throughout:
praxis \emph{executes}; rslab \emph{measures, normalizes, schemas,
receipts, and renders}; \textsf{mfact} \emph{admits formal standing}.
Each verb names a different kind of transform, and none substitutes for
another. rslab's role is specifically the middle one: it does not run
praxis-graphlaw's code and it does not decide \textsf{procint}'s
kernel-admitted status; it takes praxis's raw command output and turns it
into something a reader can trust was not hand-typed or silently edited
after the fact.

Concretely, this happens in two receipted steps. Raw command output from
a praxis run, $O$ in this thesis's $R_B \vdash A = \mu(O^*_B)$ vocabulary
(Chapter~\ref{ch:mfact}), is not itself evidence in rslab's doctrine: a
collection script, \texttt{rslab/scripts/collect\_praxis\_graphlaw.py},
schema-validates the raw output, hashes every raw file with BLAKE3, and
records the toolchain and the commit under test, producing $O^*$ ---
schema-validated, receipted output conforming to
\texttt{rslab/schemas/benchmark\_result.schema.json}. A second script,
\texttt{rslab/scripts/render\_paper\_fragments.py}, renders the LaTeX
paper fragments reported in this chapter from that receipt alone,
re-verifying every raw-file hash before rendering and refusing with
\texttt{ARTIFACT\_DRIFT\_REFUSED} on any mismatch between the receipted
hash and the file on disk. Both scripts, and the fragments they produce,
are ledgered in \texttt{.mfact/artifacts.toml} and covered by
\texttt{just regen-check}: a stale or hand-edited rslab fragment fails
the identical drift gate that a hand-edited \textsf{procint} module or
release fragment would fail (Chapter~\ref{ch:mfact}). The receipt
discipline that Chapter~\ref{ch:mfact} developed for kernel-admitted
theorems is, in this narrower sense, reused for empirical measurement:
what changes between the two rails is not whether there is a receipt, but
what the receipt is a receipt \emph{of}.

\subsection{The declared \texorpdfstring{$<$}{<} extracted \texorpdfstring{$<$}{<} stated \texorpdfstring{$<$}{<} proven ladder}
\label{sec:emp-ladder}

rslab evidence follows its own status ladder --- declared $<$ extracted
$<$ stated $<$ proven --- and \texttt{paper/main.tex} is explicit that
this ladder is \emph{deliberately never conflated} with \textsf{procint}'s
kernel-checked proven status. A \textsf{procint} declaration reaches
``proven'' only by passing kernel admission and the axiom audit
(Chapter~\ref{ch:procint}); an rslab measurement reaches the top of its
own ladder by a different route entirely, one this chapter does not
attempt to force into equivalence with the first. The praxis-graphlaw
evidence reported in this chapter sits at the \emph{extracted} rung: it
has been observed once, receipted by BLAKE3 hash and toolchain/commit
record, and is reproducible from the committed raw output --- but it is
not \emph{proven} in either rslab's own sense (which would require
further, currently absent, argument that the measurement is stable or
representative beyond the single receipted run) or in \textsf{procint}'s
kernel-checked sense (which the ladder does not even purport to reach,
since no Lean statement is at stake). Keeping the ladder and its labels
separate from \textsf{procint}'s proven/stated split is not a stylistic
choice; it is the same discipline Chapter~\ref{ch:mfact} argues for
generally, applied to a domain --- empirical software measurement ---
where the relevant admission criterion is receipted reproducibility
rather than kernel type-checking, and where collapsing the two would
silently borrow a proof-strength claim the evidence does not support.

\section{Benchmark Evidence}
\label{sec:emp-benchmark}

The praxis-graphlaw experiment, \texttt{rslab/experiments/praxis\_graphlaw},
is rslab's first instance of this pipeline. Two tables are rendered by
\texttt{rslab/scripts/render\_paper\_fragments.py} from
\texttt{rslab/receipts/praxis\_graphlaw\_benchmark\_receipt.toml}: a
conformance and cross-crate end-to-end admission-gate suite (covering
denial, SHACL-violation, N3-materialization ordering, and determinism
cases), and a benchmark suite spanning the \texttt{bencher},
\texttt{criterion}, and \texttt{divan} harnesses exercised across the
crate and the workspace root. No cell in either table is hand-typed;
both are regenerated from the same receipted raw output that
\texttt{collect\_praxis\_graphlaw.py} hashed and schema-validated.

% AUTO-RENDERED by rslab/scripts/render_paper_fragments.py — do not edit by hand.
% Regenerate: python3 rslab/scripts/collect_praxis_graphlaw.py && python3 rslab/scripts/render_paper_fragments.py
% Source: rslab/receipts/praxis_graphlaw_benchmark_receipt.toml (EXTRACTED tier;
% empirical evidence, not a Lean proof — see rslab/README.md doctrine).
\begin{table}[h!]
\centering
\begin{tabular}{lrrr}
\toprule
\textbf{Command} & \textbf{Duration (s)} & \textbf{Passed} & \textbf{Failed} \\
\midrule
\texttt{cargo test -p praxis-graphlaw} & 11.06 & 380 & 0 \\
\texttt{cargo test -p ggen --test graphlaw\_e2e} & 9.83 & 5 & 0 \\
\midrule
\textbf{Total} & & \textbf{385} & \textbf{0} \\
\bottomrule
\end{tabular}
\caption{praxis-graphlaw conformance and end-to-end admission test counts, extracted from the Ticket~018 raw command output and receipted by \texttt{rslab/scripts/collect\_praxis\_graphlaw.py} (commit \texttt{c745c14f6468}). EXTRACTED-tier empirical evidence, not a Lean proof obligation.}
\label{tab:rslab-tests}
\end{table}

\begin{table}[h!]
\centering
\begin{tabular}{lrr}
\toprule
\textbf{Command} & \textbf{Duration (s)} & \textbf{Benchmarks measured} \\
\midrule
\texttt{cargo bench -p praxis-graphlaw} (4 bench files) & 456.67 & 24 \\
\texttt{cargo bench} (3 bench files) & 379.23 & 14 \\
\midrule
\textbf{Total} & & \textbf{38} \\
\bottomrule
\end{tabular}
\caption{praxis-graphlaw benchmark suite coverage: per-benchmark output lines counted per harness (\texttt{bencher}/libtest \texttt{bench:} lines, \texttt{criterion} \texttt{time:} blocks, \texttt{divan} leaf rows); per-harness timings are not directly comparable, see receipted caveats. Per-benchmark timing detail lives in the receipted raw output, not reproduced here.}
\label{tab:rslab-benches}
\end{table}


\subsection{Receipted Caveats}
\label{sec:emp-caveats}

Three caveats travel with this evidence, receipted alongside it rather
than reported separately, and this chapter reproduces them rather than
softening them. First, profiling and flamegraph tooling did not exist in
the praxis workspace at collection time, so
\texttt{profiler\_result.schema.json} remains unpopulated; no
per-call-site cost breakdown is claimed here. Second,
``transaction-path admission control'' is not an implemented interface
in praxis-graphlaw as collected: the nearest existing structure is the
SHACL/ShEx admission gate together with \texttt{bcinr\_powl}'s admission
context described in Section~\ref{sec:emp-praxis-graphlaw}, and any
future work under the ``transaction-path'' name is a design objective
recorded for later work, not a shipped feature this chapter reports
evidence for. Third, the three benchmark harnesses report timings under
incompatible units and sampling methodologies, so this chapter, like
\texttt{paper/main.tex}, draws no cross-harness speed comparison from the
rendered table; the table reports coverage (which harnesses ran, how
many benchmarks each measured) rather than a single ranked performance
claim.

\section{Why the Rails Stay Separate}
\label{sec:emp-why-separate}

It is worth stating plainly why this chapter does not attempt to fold
praxis-graphlaw's 385 passing conformance and admission-gate cases (see
the rendered table above) into anything resembling \textsf{procint}'s
kernel-admitted theorem count. The two numbers answer different
questions under different admission criteria. A \textsf{procint} theorem
count answers: how many propositions has the Lean kernel accepted, under
a fixed and auditable axiom set, as inhabited by a checked term
(Chapter~\ref{ch:procint}, Chapter~\ref{ch:evaluation})? A
praxis-graphlaw test count answers: how many test cases, executed once
under a recorded toolchain and commit, returned the expected pass/fail
outcome, with that single run's raw output receipted so a later party
can check it was not altered after collection? The second question is a
legitimate and useful one, and rslab's receipt discipline gives an
honest answer to it. But it is not the first question, and reporting
both in the same paper --- as \texttt{paper/main.tex} does, and as this
thesis does here --- is only honest if the reader is never left to infer
that a receipted test-suite pass count carries the same epistemic weight
as a kernel-admitted proof. That is the discipline this chapter has tried
to preserve: praxis executes, rslab receipts what it measures, and
\textsf{mfact} alone admits formal standing, with the boundary between
the third activity and the first two kept visible rather than blurred by
shared vocabulary.

\chapter{AI-Assisted Formal Development}
\label{ch:ai-development}

The chapters up to this point have described a manufacturing pipeline
whose central move is treating every candidate producer as untrusted
with respect to the formal corpus, regardless of what that producer is.
This chapter turns the same discipline onto the process that actually
built the pipeline reported in this thesis: Claude Code and related
language-model tooling were used, throughout the run described in
Chapter~\ref{ch:procint} and Chapter~\ref{ch:empirical}, to generate
candidate patches, proof sketches, ontology fragments, documentation
drafts, and release metadata. This chapter asks what standing that
assistance has, what structural guardrails exist to keep it from
acquiring standing it has not earned, and what happened the one
documented time a gap in those guardrails let a claim through anyway.
It also surveys two unledgered hand-authored surfaces --
\texttt{pylab/} and \texttt{procint/Playground/} -- that exist
precisely so that AI-assisted and human experimentation has somewhere
to happen \emph{without} implicitly claiming release standing, and
closes with a case study in this project's own self-correction
discipline: a five-rail audit that found a real gap and the mechanized
scanner built in direct response to it.

\section{The Untrusted-Producer Doctrine, Applied to Its Own Construction}
\label{sec:ai-doctrine}

paper/main.tex's disclosure of generative-AI tool use states the
governing separation directly: language-model outputs generated during
the manufacturing run were treated throughout as untrusted candidate
artifacts, admitted into \textsf{procint} only when accepted by the Lean
kernel, the sorry census, the axiom audit, schema validation against the
\texttt{procint:} ontology, and the negative-fixture and release-gate
checks developed in Chapter~\ref{ch:mfact} and Chapter~\ref{ch:procint}.
No language-model output is part of the trusted base; that base is the
Lean kernel itself, Mathlib and CSLib~\cite{mathlib2020,cslib2026} at
their pinned revisions, and the \texttt{lake}/\texttt{ggen} toolchain.
The author of the accompanying paper reviewed and accepts responsibility
for all submitted content, code, citations, and claims in both the paper
and the repository -- a human accountability statement that this
chapter's discussion of automated tooling does not dilute: nothing
described below is offered as a substitute for that review, only as
machinery that makes the review's job more tractable and its omissions
more detectable.

This separation is not a precaution taken in the abstract. General-purpose
language models have a documented Lean~3/Lean~4 syntax-confusion failure
mode, generating deprecated Lean~3 constructs despite explicit Lean~4
instructions, with GPT-4-class and substantially smaller models reported
at comparable, low accuracy on the MiniF2F-Valid
benchmark~\cite{langconfusion2025}. Two families of approach trade some
of the resulting untrusted-candidate distance for a higher raw acceptance
rate: fine-tuning a general-purpose model toward Lean~4
specifically~\cite{theoremllama2024}, or pairing a language model with
solver-guided refinement~\cite{apollo2025}. \textsf{mfact} does not take
either trade in the run this thesis reports: every candidate, however
produced, passes through the same kernel-and-audit gate regardless of
its source's measured reliability, which is precisely why the framework
developed in Chapter~\ref{ch:mfact} needed to be built at all rather than
substituted for by a higher-precision generator. A more capable producer
changes how often a candidate is admitted; it does not change whether
admission, rather than production, is what confers standing.

\section{Genetic Tactic Search: Kernel-Accept Rate as the Only Fitness Signal}
\label{sec:ai-genetic-search}

One concrete, currently exploratory piece of tooling makes the
untrusted-producer doctrine's consequences for automated search
explicit. \texttt{docs/genetic-tactic-search.md} records a
genetic-programming loop, modeled on the TPOT2 pipeline-search paradigm
of evolving and selecting structured candidates by mutation and
crossover, applied here to sequences of Lean tactics rather than
scikit-learn pipeline stages. At the time that design note was written,
two declarations in \texttt{research/wfnet/obligations.toml} --- the
crown-jewel short-circuit soundness equivalence and the unfolding
correctness statement --- were both \texttt{status = "stated"}, and the
note is explicit that closing either requires new lemmas, not
recombination of existing tactics, so the search does not target them
directly. The crown-jewel obligation has since been discharged (its
current classification, \CrownJewelStatus, is rendered below); the
warm-up strategy the design note describes remains apt regardless, and
today applies to whatever open obligations remain, currently the single
declaration discussed in
Section~\ref{sec:mfact-boundary} of Chapter~\ref{ch:mfact}.

\begin{center}
% GENERATED FILE — DO NOT EDIT BY HAND (source: quadrature graph via ProcInt.crownJewel_status)
The release-time classification of the short-circuit soundness
characterization is \textbf{\CrownJewelStatus} (ground truth:
\texttt{ProcInt.crownJewel\_status}); this fragment is generated from the
catalog and cannot report it as proven while the obligation is open.

\end{center}

The search loop itself runs entirely outside Lean, in
\texttt{scripts/genetic\_tactic\_search.py}: each candidate genome is a
sequence of tactic calls drawn from a small, hand-vetted vocabulary ---
\texttt{aesop}, \texttt{simp}, \texttt{omega}, \texttt{exact},
\texttt{apply}, \texttt{rcases}, \texttt{constructor}, \texttt{intro},
\texttt{decide} --- spliced into a scratch \texttt{.lean} file and
checked with \texttt{lake env lean}. \textbf{Fitness is kernel
accept/reject, a hard, computed signal, never an LLM self-report or a
heuristic estimate}: nothing in this tool claims a proof exists without
the Lean kernel confirming it, following the same computed-not-asserted
invariant the release manifest itself obeys. A secondary, shaping term
in the fitness function draws on the empirical finding that tactics
conforming to expert proof patterns are more likely to advance a proof
successfully than arbitrary tactic
choices~\cite{tacticexpertalign2026}; this term only breaks ties among
sequences that already type-check, or biases mutation toward more
plausible candidates -- it never substitutes for the primary
kernel-accept signal. This is a modest but real precedent: native
LLM-tactic integration into Lean~4 has reached substantial per-step
automation rates elsewhere (Lean Copilot reports 74.2\% step-level
automation against a 40.1\% \texttt{aesop}-only
baseline~\cite{leancopilot2024}), which motivates treating expert-pattern
conformance as plausible signal, without licensing the tool to call an
LLM directly for tactic proposal in its first version -- a restraint the
design note attributes in part to the syntax-confusion risk already
discussed in Section~\ref{sec:ai-doctrine}~\cite{langconfusion2025}.
Closing genuinely new theorems, as opposed to routine proof steps,
remains hard even for purpose-built systems: HyperTree Proof Search
reports 82.6\% on Metamath with online
training~\cite{hypertree2022}, and HOList's benchmark targets a large but
bounded tactic library rather than open-ended lemma
discovery~\cite{holist2019}. Both results are cited in the design note as
the reason the tool's v1 target is warm-up lemmas in
\texttt{procint/ProcInt/Playground/TacticSearchWarmup.lean}, not either
crown-jewel-class obligation directly. Infrastructure the design note
surveys as relevant context but does not depend on in v1 includes
Lean4Lean, an independently verified Lean~4
typechecker~\cite{lean4lean2024}; Lean-auto, an interface to external
automated theorem provers~\cite{leanauto2025}; LeanExplore, a semantic
and lexical search engine over Lean declarations useful for premise
selection if the tool grows beyond a fixed vocabulary~\cite{leanexplore2025};
and large synthetic Mathlib4-derived theorem corpora for training
proof-synthesis models at scale~\cite{mathlib4dataset2025}.

The candidate-to-admission boundary for this tool is stated as
precisely as the rest of the pipeline demands. A kernel-accepted tactic
sequence the script finds is a legitimate \texttt{Mfact.Candidate} with
\texttt{producer = "genetic\_tactic\_search"}, in the sense of
Section~\ref{sec:mfact-candidates}. Promotion to \texttt{Admitted} --
flipping a TTL fragment's \texttt{status} from \texttt{"stated"} to
\texttt{"proven"} and re-running the full \texttt{just render \&\& just
build \&\& just audit \&\& just manifest \&\& just certify} chain -- is a
human-reviewed, separate step; the script itself never writes to
\texttt{packs/*/fragments/*.ttl}. Success for the search is defined
narrowly, as \texttt{lake env lean <candidate-file>} exiting~0, logged
verbatim; the design note is explicit that no claim of a
\textsc{Stated}-to-\textsc{Proven} transition is made without the full
\texttt{just check} / \texttt{just manifest \&\& just certify} chain
passing on a candidate a human has promoted. The tool's own status
line is honest about its place in the larger system: exploratory,
off-ledger tooling, not part of the certified pipeline.

\section{Unledgered Hand-Authored Surfaces: \texttt{pylab/} and \texttt{procint/Playground/}}
\label{sec:ai-unledgered}

Section~\ref{sec:mfact-ledger} of Chapter~\ref{ch:mfact} introduced the
ledger law's two permanently excluded surfaces by name;
this section develops why they exist and what governs them.
\texttt{procint/Playground/} is the pre-existing hand-authored, unledgered
Lean demonstration surface: worked, runnable instances of algorithms the
repository cares about, proving that \textsf{procint}'s own formal
content is actually exercised, not merely declared, without those
demonstrations themselves acquiring release standing. \texttt{pylab/} is
the Python-language counterpart, existing for the same reason but for
tools that have no reasonable Lean formalization: AutoML pipeline
search, the \texttt{pm4py} process-mining ecosystem, and LLM-to-PDDL
generation. Both surfaces share a governance tier, differing only in
language: neither is ggen-rendered, neither is ledgered in
\texttt{.mfact/artifacts.toml}, and neither carries standing, counts, or
certification data -- their absence from the ledger is, per AGENTS.md,
correct by design, not an omission awaiting correction. \texttt{pylab/}
is explicitly never added to \texttt{just build}, \texttt{just check},
\texttt{just release}, or the artifact ledger; a failure inside it is an
ordinary code-review issue, never \texttt{ARTIFACT\_DRIFT\_REFUSED} or a
release-standing regression, and PYLAB-PRD-ARD-v26.7.7.md is explicit
that nothing in \texttt{pylab/} is claimed to be semantically equivalent
to, or a proof about, any Lean formalization in
\texttt{procint/ProcInt/} -- the same restraint the repository applies to
its own \texttt{ProcInt.Planning.Pddl} module, whose relationship to any
future \texttt{pddl-plus-parser}-based Python experiment is, in the PRD
document's own words, ``two independent things that happen to model
similar concepts,'' not a proven-equivalent pair, absent an actual
correspondence proof this repository does not currently attempt.

\texttt{pylab/}'s tool inventory, researched and wired via \texttt{uv
add} rather than direct \texttt{pyproject.toml} edits, comprises five
research libraries plus the FastMCP server framework used to expose
them. \textbf{TPOT2} (merged into mainline \texttt{tpot} v1.1.0)
provides genetic-programming AutoML, evolving scikit-learn pipelines as
DAGs under multi-objective selection, and is the direct model for the
tactic-search loop of Section~\ref{sec:ai-genetic-search}. \textbf{pm4py}
is the foundation process-mining library and the confirmed reference
implementation of both OCEL~2.0 and Improved Token-Based Replay, serving
as a dependency for the two libraries that follow it. \textbf{powl}
provides POWL~v2 discovery and Petri-net/BPMN/POWL conversion, hard-
depending on \texttt{pm4py}; at the time of the PRD document, a real
version incompatibility existed between \texttt{powl==2.3.7} and
\texttt{pm4py==2.2.32} (a missing internal \texttt{pm4py} module path
\texttt{powl} expects), handled in \texttt{pylab/}'s wrapper by a lazy
import and a structured error rather than a crash. \textbf{ocpa}
supplies object-centric process analysis over OCEL~2.0 data, including
sqlite import and object-centric Petri-net discovery. \textbf{pddl-plus-parser}
parses PDDL+, a continuous/hybrid superset of the PDDL~3.1 fragment
already covered by the repository's Rust \texttt{pddl} crate, and is
retained as a genuinely complementary parser for benchmark sourcing and
cross-validation rather than a duplicate.

Two decisions recorded in the PRD document are worth treating as
governance case studies in their own right, because both concern tools
that were seriously considered and then not adopted, for reasons the
document is careful to record so the decision is not re-litigated.
\textbf{SpiffWorkflow} was considered as a live POWL~v2 execution engine
and explicitly rejected: it is fundamentally BPMN-token-based, compiling
sequence flows and gateway types into a \texttt{TaskSpec} tree, and
feeding it a POWL partial-order block would require a lossy graph
transformation, since arbitrary partial orders do not reduce cleanly to
nested AND/XOR gateways without auxiliary dummy tasks. Since the
sibling \texttt{bcinr-powl} crate already implements POWL~v2 execution
natively with an exact op-kind semantics, there was no reason to adopt a
lossy BPMN intermediary in its place. \textbf{l2p}, an LLM-to-PDDL tool
generating domain and problem PDDL from natural language, was scoped in
the PRD document's first revision as a candidate front end
(\texttt{breed\_l2p}, with a planned FastMCP wrapper); the revision
actually applied instead built a general-purpose FastMCP server exposing
read-only Lean/Lake/\texttt{just} introspection and one tool per
\texttt{pylab} library, superseding the l2p-specific wrapper. AGENTS.md's
governing rule since goes further than that supersession: \texttt{l2p} is
to be deprecated outright -- not installed, referenced, or used as a
planning library -- a stricter position than the PRD document's own
``decided: built with FastMCP'' framing for it. This chapter treats the
later, stricter AGENTS.md rule as authoritative and the PRD document's
earlier framing as superseded history, which is itself a small instance
of the governance discipline this chapter is about: a decision recorded
once is not immutable, and a later, more restrictive rule in AGENTS.md
governs current and future work regardless of what an earlier planning
document scoped.

\section{The Standing-Guard Scanner: Read-Only Gap Detection}
\label{sec:ai-standing-guard}

\texttt{pylab/src/mpops/standing\_guard/server.py} is itself an instance
of the \texttt{pylab/} surface just described: an unledgered,
hand-authored FastMCP server, described in its own module docstring as
``read-only integrity, correctness, and guardrail scan check tools.'' It
exposes a single tool, \texttt{scan}, which runs eight independent gap
checks over the repository and returns each finding as a typed record
carrying a gap class, a severity, a refusal code, the path or target
affected, observed evidence, an expected value, an actual value, and a
recommended action -- itself a small, non-Lean echo of the typed-refusal
discipline of Section~\ref{sec:mfact-candidates}. The eight checks, in
the order the scanner runs them, are:

\begin{enumerate}
\item \textbf{Sorry theorem promotion.} The scanner parses every TTL
  fragment under \texttt{packs/lean-math-pack/fragments/} plus
  \texttt{ontology.ttl} for declarations whose \texttt{procint:status} is
  \texttt{"proven"}, compiles a single scratch Lean file importing each
  declaration's module and running \texttt{\#print axioms} on each
  proven name via \texttt{lake env lean --stdin}, and flags
  \texttt{SORRY\_THEOREM\_PROMOTED} if the kernel output for that name
  mentions \texttt{sorryAx}, or if the name is unknown to the Lean
  environment at all.
\item \textbf{Ledger drift.} The scanner walks every artifact recorded
  in \texttt{.mfact/artifacts.toml}, recomputes a BLAKE3 hash of the file
  on disk (falling back to the \texttt{b3sum} CLI if a native hash
  implementation is unavailable), and flags \texttt{ARTIFACT\_MISSING} if
  the file is absent, \texttt{ARTIFACT\_DRIFT\_REFUSED} if the computed
  hash disagrees with the ledgered hash, and \texttt{ORPHAN\_ARTIFACT\_REFUSED}
  if \texttt{git ls-files} or \texttt{git status --porcelain} reports the
  ledgered file as untracked.
\item \textbf{Orphan artifact scan.} The scanner globs
  \texttt{release/**/*.json} and \texttt{paper/**/*.tex} (excluding
  \texttt{procint/Playground/**}, \texttt{pylab/**}, and
  \texttt{paper/main.tex} itself) for files absent from the ledger, and
  flags \texttt{ORPHAN\_ARTIFACT\_REFUSED} when such a file's content
  looks standing-bearing -- containing any of \emph{proven}, \emph{stated},
  \emph{certified}, \emph{sorry}, \emph{hash}, \emph{audit},
  \emph{quadrature}, or \emph{verif} -- or when it already lives under
  \texttt{release/} or \texttt{paper/} regardless of content.
\item \textbf{regen-check coverage gap.} The scanner extracts the body of
  the justfile's \texttt{regen-check:} recipe and checks, for each
  ledgered artifact, whether its declared producer -- \texttt{ggen}, a
  script path, or another named command -- is textually referenced
  inside that recipe body, flagging \texttt{REGEN\_CHECK\_COVERAGE\_GAP}
  as a warning when it is not: a producer that only runs under some
  other, separately invoked recipe is a coverage gap in
  \texttt{regen-check}, not a convenience.
\item \textbf{Correspondence binding.} The scanner reads
  \texttt{release/verif-receipt.json} and, for each recorded obligation,
  flags \texttt{STALE\_PROOF\_BINDING} if its \texttt{aeneasDecl} field is
  the literal string \texttt{"TBD"}, or if its status is
  \texttt{PROVEN} but the corresponding Lean proof file under
  \texttt{dist/verif/lean/.../Corr/} is missing, or its comment-stripped
  body does not reference the declared \texttt{aeneasModule}. This check
  is developed further, alongside the correspondence rail proper, in
  Chapter~\ref{ch:correspondence}.
\item \textbf{Tag ancestry.} The scanner runs \texttt{git rev-parse} and
  \texttt{git merge-base --is-ancestor} to confirm that the tag
  \texttt{v26.7.7-procint-certified} is an ancestor of \texttt{HEAD} and,
  when \texttt{release-manifest.json} carries a \texttt{runIdentifier},
  that the identifier is itself an ancestor of the tag, flagging
  \texttt{TAG\_ANCESTRY\_FAIL} otherwise -- the mechanized instance of the
  core-identity protection developed in
  Section~\ref{sec:mfact-core-identity}.
\item \textbf{Untracked ontology fragments.} The scanner globs
  \texttt{packs/*/fragments/*.ttl} and flags
  \texttt{UNTRACKED\_ONTOLOGY\_FRAGMENT} for any fragment file git
  reports as untracked, closing the gap between ``committed source
  tree'' and ``what \texttt{cat fragments/*.ttl} actually concatenates.''
\item \textbf{Prose/paper consistency.} The scanner reads
  \texttt{paper/main.tex} line by line and paragraph by paragraph. Two
  rules are self-contained per line: a stale bare count of 145 or 318
  (the pre-\TotalDecls totals reported in STANDING.md's own historical
  ladder) flags \texttt{STALE\_PAPER\_PROSE\_COUNT}, and a phrase
  matching ``Aeneas proves/verified/checked/certified'' flags a
  \texttt{PROSE\_LINT\_VIOLATION} in favor of the more precise ``Aeneas
  extracts'' and ``Lean proves.'' Three further rules are scoped to the
  whole paragraph rather than a single line, to avoid false positives
  from LaTeX's hand-wrapped line breaks splitting a trigger word from its
  qualifier: a bare \emph{chain} or \emph{hash} must be accompanied
  somewhere in the same paragraph by a qualifying phrase such as
  ``receipt chain'' or ``content hash''; a bare \emph{prove}/\emph{proof}/
  \emph{proven} must be accompanied by a formal-context word such as
  \emph{Lean}, \emph{kernel}, or \emph{axiom}; and a totality adverb --
  \emph{entirely}, \emph{completely}, \emph{fully}, \emph{absolutely},
  \emph{always}, \emph{never} -- must be accompanied by an explicit scope
  qualifier. Each unaccompanied occurrence is reported as a
  \texttt{PROSE\_LINT\_VIOLATION} warning at the specific line the
  trigger word occurs on.
\end{enumerate}

The scanner is deliberately read-only: it reports findings, it does not
repair the repository, and its own existence inside \texttt{pylab/} means
it is subject to exactly the same governance tier -- no standing claim,
no ledger entry -- as the surface it audits. It is, in effect, the
Section~\ref{sec:ai-doctrine} doctrine applied reflexively: automated
tooling checking the pipeline's own discipline is itself treated as an
untrusted, unledgered producer of findings, not as a certifying
authority in its own right.

\section{The 2026-07-07 Five-Rail Audit: A Case Study in Self-Correction}
\label{sec:ai-guardrail-audit}

The standing-guard scanner's eight checks did not arise from abstract
risk analysis. A five-rail audit conducted on 2026-07-07, recorded in
\texttt{pylab/docs/jira/26.7.7/tickets/ticket\_013\_v26\_7\_7\_gap\_audit.md},
found that a sorry-backed theorem had been ledgered with
\texttt{procint:status "proven"} in the TTL catalog, with no guard in
place that would have caught it, alongside several related silent-drift
failure modes. AGENTS.md records the resulting structural rules under a
``Guardrails (post-v26.7.7-audit)'' heading, framed explicitly as making
each specific failure mode \emph{structurally impossible}, not merely
discouraged, and each rule maps onto one of the scanner checks of
Section~\ref{sec:ai-standing-guard}:

\begin{itemize}
\item \textbf{No status promotion without a checked guard.} Any TTL
  fragment setting \texttt{procint:status "proven"} on a theorem must
  have a corresponding entry in \texttt{release/gates.json}, or an
  equivalent builder check, that mechanically verifies \texttt{\#print
  axioms} shows no \texttt{sorryAx} for that declaration; a status
  literal in TTL alone is never sufficient, and a human or agent setting
  \texttt{"proven"} without a passing mechanical check is
  \texttt{STATED\_PROMOTED\_TO\_PROVEN}. This rule's mechanized instance
  is the scanner's sorry-theorem-promotion check
  (Section~\ref{sec:ai-standing-guard}, item~1). A named companion
  vocabulary applies specifically to the five deliberately stated
  workflow-countermodel declarations introduced in
  Section~\ref{sec:mfact-boundary}: their status key is
  \texttt{WFNET\_INFINITE\_TRANSITION\_COUNTERMODEL}, the guard that must
  never let them be promoted is named \texttt{countermodel\_not\_promoted},
  and the refusal that guard fires is \texttt{COUNTERMODEL\_PROMOTION\_REFUSED}
  -- language that exists precisely so these five declarations are never
  mistaken for unfinished work awaiting the same promotion path as an
  ordinary stated theorem.
\item \textbf{regen-check must not have untracked-file blind spots.}
  \texttt{git diff --exit-code} only sees tracked files, so a ledgered
  artifact's producer script must itself be invoked by
  \texttt{regen-check}/\texttt{check}, not merely by some separate recipe
  outside that chain; newly created ledgered artifacts must be
  \texttt{git add}ed in the same commit that introduces them, and an
  untracked ledgered file whose disk hash disagrees with
  \texttt{.mfact/artifacts.toml} is \texttt{ORPHAN\_ARTIFACT\_REFUSED} and
  must actually fire, not merely exist as a documented category. This
  rule's mechanized instances are the scanner's ledger-drift and
  regen-check-coverage-gap checks (Section~\ref{sec:ai-standing-guard},
  items~2 and~4).
\item \textbf{Correspondence theorems must reference their extraction.}
  A correspondence obligation reported as \texttt{PROVEN} must have its
  Lean statement actually import and quantify over the extracted or
  generated declaration -- for example an Aeneas \texttt{Generated.*}
  type -- via its declared abstraction function, not merely present a
  same-shaped statement over hand-written types; a receipt field such as
  \texttt{aeneasDecl: "TBD"} is itself a refusal signal that must block
  \texttt{PROVEN} status rather than ship alongside it. This rule's
  mechanized instance is the scanner's correspondence-binding check
  (Section~\ref{sec:ai-standing-guard}, item~5), developed further in
  Chapter~\ref{ch:correspondence}.
\item \textbf{Tag ancestry is part of certification.}
  \texttt{just certify}/\texttt{just release} must verify the release
  commit is an ancestor of, or equal to, the certified git tag before
  reporting \texttt{CERTIFIED\_RELEASE=PASS}; a tag that predates the
  currently rendered artifacts is \texttt{RECEIPT\_RECURSION\_REFUSED},
  not a historical curiosity to note and move past. This rule's
  mechanized instance is the scanner's tag-ancestry check
  (Section~\ref{sec:ai-standing-guard}, item~6), and its target property
  is exactly the core-identity protection of
  Section~\ref{sec:mfact-core-identity}.
\item \textbf{New fragments feeding \texttt{ontology.ttl} must be
  tracked before \texttt{regen-check} runs.} An untracked \texttt{.ttl}
  file silently concatenated into generated output is a source-provenance
  gap distinct from ordinary generated-artifact drift: it means the
  generated output is not reproducible from the committed source tree at
  all. This rule's mechanized instance is the scanner's untracked-fragment
  check (Section~\ref{sec:ai-standing-guard}, item~7).
\end{itemize}

Read together, the audit and its resulting guardrails are the strongest
evidence this thesis can offer, internal to the repository itself, that
the untrusted-producer doctrine of Section~\ref{sec:ai-doctrine} is
enforced rather than merely stated. The failure the audit found was
real, not hypothetical: a status literal reached \texttt{"proven"} in a
committed fragment without the mechanical check that should have gated
it. The response was not a promise to be more careful, but a scanner
whose sorry-theorem-promotion check, run against the repository today,
would have caught exactly that failure, together with six further checks
covering the related silent-drift modes the same audit identified. This
chapter's account of AI-assisted development therefore closes where
Chapter~\ref{ch:mfact} opened: standing is never inferred from the fact
that a candidate -- human-authored, agent-authored, or produced by a
genetic search over tactic sequences -- made contact with the
repository. It is conferred only by an admission step that a later,
independent check can re-run and, when it fails, is designed to catch.

\section{Summary}

This chapter examined the pipeline's own construction under the same
discipline the pipeline enforces on \textsf{procint}'s mathematical
content. Generative-AI tooling is treated, by explicit disclosure and by
mechanism, as an untrusted candidate producer whose output enters the
trusted base only through kernel admission, sorry census, axiom audit,
and release-gate checks -- never through the fact of having been
produced by a more or less capable model. The genetic tactic search tool
instantiates this doctrine for automated proof search specifically:
fitness is kernel accept/reject, never an LLM's self-report, and
promotion from a found candidate to an admitted theorem remains a
separate, human-reviewed step outside the search script's own reach.
\texttt{pylab/} and \texttt{procint/Playground/} provide governed spaces
for hand-authored experimentation -- including the tool inventory and
the SpiffWorkflow and l2p decisions recorded above -- that never acquire
release standing by virtue of existing. And the standing-guard scanner,
together with the five-rail audit and its resulting AGENTS.md guardrails,
is this project's own record of a real gap in that discipline and the
mechanized, read-only check subsequently built to close it. Chapter~\ref{ch:evaluation}
now turns to consolidating the quantitative release evidence this and
earlier chapters have referenced piecemeal into a single evaluation
treatment.

\chapter{Evaluation}
\label{ch:evaluation}

Earlier chapters introduced the release's quantitative evidence
piecemeal, in service of whatever argument was at hand: Chapter~\ref{ch:procint}
cited the declaration catalog's proven/stated split while presenting
\textsf{procint}'s module architecture; Chapter~\ref{ch:correspondence}
cited the correspondence rail's status while presenting the D1 obligation;
Chapter~\ref{ch:empirical} cited rslab's receipted test and benchmark
counts while keeping them apart from kernel-checked standing. This
chapter consolidates that evidence into a single evaluation treatment,
following the structure of \texttt{paper/main.tex}'s own Evaluation
section: it explains what each rendered table measures and how the
number was computed, then defers to the rendered fragment itself for the
number, rather than restating any count, hash, or status literal by hand
in this chapter's prose. As Chapter~\ref{ch:mfact} argued and
Chapter~\ref{ch:procint} exercised, a number typed into prose is exactly
the kind of unreceipted assertion this thesis's manufacturing discipline
exists to avoid; every quantitative claim below is either a macro from
\texttt{release\_macros.tex} or an \verb|\input| of a ledgered fragment
rendered from \texttt{release-manifest.json} or \texttt{quadrature.json}.

\section{Formal Standing Evaluation}
\label{sec:eval-formal}

The first table this chapter reports is rendered directly from the
release manifest of the certified run, \texttt{release-manifest.json},
and from nothing else: no cell is written by hand, and the generation
pipeline fails closed, so that if the manifest is absent or stale the
fragment does not render and the surrounding document does not build for
release. What the table measures is the release's own accounting of its
declaration catalog: how many of the \TotalDecls\ declarations recorded
at release time reached kernel-admitted, axiom-audited proven status
(\KernelTheorems), how many remain formally stated with an open proof
obligation (\StatedCount), and the release-wide \texttt{sorry} and
unauthorized-\texttt{admit} counts (\SorryCount\ and \AdmitCount\
respectively, both required to be zero for \CertifiedStatus\ to read
PASS). These are not independently re-derived by this thesis; they are
the same manifest values \texttt{release\_macros.tex} exposes as
\LaTeX\ macros and Chapter~\ref{ch:procint} has already cited in context.

\input{../paper/evaluation}

\paragraph{Final status.} Table~\ref{tab:mfact-final-status} in
Chapter~\ref{ch:mfact}, generated from the release manifest and the
standing records that accompany it, aggregates the individual gate
outcomes --- axiom audit, negative fixtures, certification, and standing
quadrature --- into the single-row summary that \texttt{just status}
(Chapter~\ref{ch:ai-development}) reads back at the command line; it is
reproduced there rather than repeated here.

\section{Standing Quadrature}
\label{sec:eval-quadrature}

The release is not evaluated only by whether the Lean packages build.
\texttt{paper/main.tex} is explicit that it is evaluated by whether the
declaration catalog, the admitted Lean corpus, the process evidence, the
release manifest, and the paper's own claims form a \emph{closed standing
graph}: the Standing Quadrature check rejects orphan declarations,
unsupported claims, untraced artifacts, and any divergence between the
manifest and the paper. This closure property is itself kernel-checked
rather than asserted by a separate script running outside Lean's trust
boundary: a generated witness module, \texttt{ProcInt.Release.Quadrature},
carries the name-sorted proven surfaces of the catalog, the manifest, and
the axiom audit as literal lists and proves their pairwise equality by
\texttt{rfl} --- so that any orphan declaration on any one side makes the
corresponding lists differ syntactically and the kernel refuses the
witness outright, rather than the check silently passing on a mismatch.
The same witness module encodes the manufacturing run itself as a
\textsf{procint} process trace and proves, using the corpus's own
\texttt{DeclareConstraint} semantics from Chapter~\ref{ch:procint}, that
the run conforms to its declared process model: \textsf{procint} is, in
this one specific and kernel-checked sense, used to model the process by
which \textsf{mfact} manufactured \textsf{procint}. The table below is
rendered from \texttt{quadrature.json}; three negative controls --- a
corrupted evaluation number, a claim with no evidence, and a catalog
declaration with no rendered symbol --- are each engineered to cause a
typed refusal rather than a silent pass, in the same spirit as the
negative fixtures Chapter~\ref{ch:procint} and Chapter~\ref{ch:mfact}
describe for the Lean corpus itself.

% RENDERED by ggen sync from the quadrature graph (quadrature-pack).
% Do not edit: regenerate via `just standing-quadrature`.
\begin{table}[h]
\centering
\small
\begin{tabular}{lrl}
\toprule
Surface pair & Count & Result \\
\midrule
TTL to Lean & 160 & PASS \\
Lean to Audit & 160 & PASS \\
Audit to Manifest & 160 & PASS \\
Manifest to Paper & 6 & PASS \\
Artifact to Process Event & 215 & PASS \\
Claim to Evidence & 5 & PASS \\
\bottomrule
\end{tabular}
\caption{Standing Quadrature: PASS over
160 proven declarations, 5 traced paper
claims, and 6 evaluation numbers. The closure is
kernel-checked by the rendered witness
\texttt{ProcInt.Release.Quadrature}.}
\label{tab:quadrature}
\end{table}


\section{Correspondence Status}
\label{sec:eval-correspondence}

Chapter~\ref{ch:correspondence} developed the Aeneas/Charon correspondence
rail in detail, including the D1 token-replay-counts obligation, the
guardrail requiring a correspondence theorem to quantify over its actual
extracted declaration rather than a same-shaped hand-written stand-in, and
that rail's own generated evaluation table (rendered by the correspondence
builder from its receipt --- a charon/aeneas pipeline receipt, a
\texttt{lake build} of the extracted module against \textsf{procint}, and
an axiom-print check on the resulting theorem --- never written by hand),
reproduced there rather than repeated here.

\section{Falsifier and Valid Objection Surface}
\label{sec:eval-falsifier}

The evaluation above answers ``what does the release currently claim.''
\texttt{paper/main.tex}'s Falsifier and Valid Objection Surface section
answers a complementary question: what would it take for one of those
claims to be wrong, and is that possibility engineered into the release
gate rather than left as an informal hope. A valid objection to the
release, in the sense that section defines, must exhibit at least one of
eight closed failure classes: a declared admitted theorem contains
\texttt{sorry}; a theorem depends on an axiom outside the authorized set;
a manifest hash fails to match its artifact; a negative fixture passes
when it is designed to fail; a generated evaluation number fails to match
the manifest that is supposed to have produced it; a formally stated
declaration is described, anywhere in the paper or its fragments, as
proven; the catalog-to-Lean rendering step changes the theorem statement
it was supposed to project; or the release cannot be replayed under its
pinned dependencies.

These eight classes are deliberately enumerated as \emph{closed} failure
classes rather than left as an open-ended attack surface: each one is
engineered into a specific gate, a typed refusal, or a non-admitted
state, so that the release gate's design goal is that any occurrence of
one of these cases prevents admission or produces a typed refusal, rather
than passing silently through to a certified release. After certification,
no constructor of this valid-objection surface is inhabited within the
declared boundary --- the surface is empty in the sense that its
negative controls, described alongside Section~\ref{sec:eval-quadrature}'s
quadrature controls and Chapter~\ref{ch:procint}'s negative fixtures,
each demonstrate the corresponding gate actually refusing on a poisoned
copy, rather than the emptiness being asserted without a check that could
have failed. It is worth restating the scope of this claim precisely, in
the same register \texttt{paper/main.tex} insists on: a local objection
to one proof, one citation, or one template is not automatically a
refutation of the release's mathematical-manufacturing claim as a whole.
It becomes one only if it crosses the admission boundary and invalidates
a specific, claimed receipt --- and the eight classes above are exactly
the enumerated ways of doing that, not a suggestive sample of a larger,
unstated set. This same enumeration is the basis for
Chapter~\ref{ch:discussion}'s discussion of what admission does, and does
not, check.

\chapter{Discussion, Limitations, and Future Work}
\label{ch:discussion}

The preceding chapters have made a sustained argument that standing must
be earned through admission rather than asserted by prose, and
Chapter~\ref{ch:evaluation} closed by enumerating the eight closed
failure classes that would falsify the release if any were inhabited.
This chapter turns that same scrutiny on the release's own boundaries:
what the proven/stated discipline does and does not claim, what kernel
admission does and does not check, what this thesis explicitly excludes,
and --- in the Future Work section --- what remains open, including a
frank account of a toolchain limitation specific to the sandbox this
thesis was written in, kept carefully separate from any claim about the
certified release itself.

\section{Proven versus Stated}
\label{sec:disc-proven-stated}

\texttt{paper/main.tex}'s Limitations and Standing section states the
library's discipline plainly: the library's standing is exactly the
manifest's proven/stated split, and prose adds nothing to it. A
\emph{stated} declaration, in this vocabulary, is a formal statement
whose proof obligation is open: it compiles under the pinned toolchain,
it is citable as a definition of the property it names, and it carries no
more authority than that. Chapter~\ref{ch:procint} exercised this
distinction on the \TotalDecls-declaration catalog, and it is worth
restating the current split precisely, without smoothing the two
categories together as this thesis's discipline forbids: \KernelTheorems\
declarations are kernel-admitted, axiom-audited proven theorems, and
\StatedCount\ remain formally stated with an open obligation.

Among those \StatedCount\ stated declarations, two categories are worth
distinguishing again here, because collapsing them is exactly the kind of
mistake the proven/stated discipline exists to prevent. The first is a
single declaration, \texttt{ProcInt.BranchingProcess.isUnfoldingOf\_statement}
(\texttt{packs/lean-math-pack/fragments/petriext.ttl}), the unfolding-correctness
statement tracked in \texttt{research/wfnet/obligations.toml} as
obligation order 3, \texttt{name = "unfolding\_correctness"},
\texttt{status = "STATED"} --- this is the one declaration in the current
catalog that represents genuine open proof work, in the sense that
closing it would add a new kernel-admitted theorem rather than merely
documenting an intentional boundary. The second category is the five
WF-net countermodel declarations in
\texttt{packs/lean-math-pack/fragments/workflow\_countermodel.ttl} ---
\texttt{crownCounter\_reaches\_mid}, \texttt{crownCounter\_reaches\_final},
\texttt{crownCounter\_sound}, \texttt{crownCounter\_not\_bounded}, and
\texttt{WfNet.infinite\_transition\_countermodel\_sound\_not\_bounded} ---
which are \emph{deliberately and permanently} stated. These five
declarations do not represent unfinished work; they witness why the
crown-jewel equivalence requires its \texttt{[Finite~T]} hypothesis in
the first place, by exhibiting a net that is sound and not bounded once
that hypothesis is dropped. A guard, \texttt{countermodel\_not\_promoted},
is engineered specifically to refuse (\texttt{COUNTERMODEL\_PROMOTION\_REFUSED})
any attempt to promote these five to proven standing, precisely because
promoting a deliberate negative witness to proven status would corrupt
the very fact it is meant to record. \texttt{just status} in this
sandbox's environment reports this family under its own status key,
\texttt{WFNET\_INFINITE\_TRANSITION\_COUNTERMODEL=STATED}, distinct from
the ordinary stated/proven split, for exactly this reason: it is a
status that is correct and intended to remain STATED, not a gap awaiting
closure alongside the unfolding-correctness obligation.

\section{Crown-Jewel Scope}
\label{sec:disc-crown-jewel}

The crown-jewel theorem --- the workflow-net soundness equivalence
(van der Aalst 1997, Lemma 8)~\cite{vanderaalst1997}, characterizing
soundness as equivalent to short-circuit liveness and boundedness under
\texttt{[Finite~T]} --- carries \CrownJewelStatus\ status at the current
release. \texttt{paper/main.tex} is careful to scope this claim per
direction: if only one direction of the equivalence were admitted at a
given release, the claim would be that direction and nothing more, with
the converse remaining a stated obligation in the manifest rather than an
implicit part of the headline result. This chapter reproduces that
caution rather than restating a stronger claim than the manifest
supports; the crown-jewel status fragment rendered in
Chapter~\ref{ch:evaluation} and Chapter~\ref{ch:procint} is the authority
on which direction, or directions, the current release actually admits.

\section{What Admission Does Not Check}
\label{sec:disc-admission-limits}

Kernel admission together with the axiom audit establishes a specific
and narrow thing: that an admitted term proves its stated theorem under
the pinned, authorized axiom set. It does not establish that the
catalog-curated statement is a faithful rendering of the informal
mathematical canon it claims to formalize. That correspondence --- does
this Lean proposition actually say what van der Aalst's Lemma~8 says, or
what the process-mining literature means by ``sound'' --- rests on
citation-annotated statement curation reviewed against primary sources,
and is exactly the kind of residual risk the semantic audits of Meek et
al.~\cite{meek2026formalizing} target: compilation success can coexist
with a silently weakened hypothesis or an incompletely stated multi-part
theorem, and no amount of kernel-side checking distinguishes a faithful
formalization from a subtly narrower one that still type-checks.
Chapter~\ref{ch:background} introduced this diagnosis; this chapter
returns to it as a standing limitation rather than a risk this pipeline
claims to have retired. Fidelity of statements is, in this repository's
own words, a review obligation, not a theorem.

A second limit concerns novelty. Any ``first mechanization'' claim made
elsewhere in this thesis or in \texttt{paper/main.tex} is scoped to a
search of the Lean (Mathlib/CSLib), Coq, and Isabelle (AFP) ecosystems
conducted at run time; it is a negative search result at a point in time,
not a proof of absence, and later work in any of those ecosystems could
in principle supersede it without this thesis's claim having been false
when made.

A third limit is that standing is not to be inferred from this document's
prose at all. At release time, the crown-jewel rail, the countermodel
rail, the correspondence rail, the product/DX rail, and the paper rail
each carry independent manifest status; the paper, and by extension this
thesis, is a publication surface reporting that status, not itself a
source of truth. Where paper prose and manifest status diverge, the
manifest and its receipts govern, and this chapter treats that as a
standing methodological commitment rather than a hypothetical.

\section{Exclusions}
\label{sec:disc-exclusions}

\texttt{paper/main.tex} enumerates, in tabular form, a set of claims this
work explicitly does not make, and this thesis inherits that enumeration
rather than relaxing it: that LLMs prove mathematics autonomously (they
are treated throughout as untrusted candidate producers, per
Chapter~\ref{ch:mfact} and Chapter~\ref{ch:ai-development}); that
\textsf{procint} formalizes all of process intelligence (the v1 boundary
is manifest-scoped, per Chapter~\ref{ch:procint}); that kernel admission
proves informal intent (statement fidelity is a review obligation, per
Section~\ref{sec:disc-admission-limits} above); that Rice's theorem is
violated (arbitrary semantic correctness is not decided by anything in
this pipeline); that external opinion enters the standing calculus (only
receipted artifacts are parsed by the certification gate); that criticism
outside the declared boundary is refuted by the empty-objection-surface
claim of Chapter~\ref{ch:evaluation} (that claim is scoped strictly to
the declared boundary and says nothing about objections outside it); that
\textsf{ggen} output has standing merely by having been generated (only
admitted output has standing, per the projection-versus-admission
distinction of Chapter~\ref{ch:mfact}); and that a stated theorem is
proven (stated means, without qualification, that the obligation remains
open).

\section{Future Work}
\label{sec:disc-future-work}

Four items make up this thesis's honest account of what remains open,
ranging from a specific closable proof obligation to a limitation of this
particular writing environment that this chapter is careful not to
conflate with the certified release's own standing.

\paragraph{(a) Closing the unfolding-correctness obligation.}
The single genuinely open item in the current \TotalDecls-declaration
catalog, \texttt{ProcInt.BranchingProcess.isUnfoldingOf\_statement}, is a
concrete, scoped target for future proof work: its statement, as recorded
in \texttt{packs/lean-math-pack/fragments/petriext.ttl}, characterizes
when a branching process is a legal unfolding of a net~$N$ by conjoining
acyclicity of the occurrence net, absence of backward conflict, and two
conditions relating the occurrence net's pre/post-condition structure to
$N$'s transitions on singleton arcs. The declaration's own doc-comment
attributes the underlying homomorphism condition it specializes to
Engelfriet 1991, Definition 3.1; this thesis reports that attribution
exactly as the source states it, in prose rather than as a citation
key, because no corresponding entry exists in \texttt{paper/refs.bib} and
this thesis's discipline forbids inventing one. Closing this obligation
means discharging its stated proposition under the same kernel-admission
and axiom-audit gate every other \KernelTheorems-count theorem in the
catalog passed through --- not weakening the statement to make it easier,
and not asserting proven status without the mechanical check.

\paragraph{(b) Un-run docs and post-release lanes.}
\texttt{just next}, run read-only in this sandbox, lists two lanes that
have not yet been attempted in this environment: the publication packet
has not been manufactured (actionable via \texttt{just manufacture-post-release}),
and the docs lane has not been attempted (actionable via \texttt{just docs}).
Neither omission changes the certified release's proven/stated split or
gate status; they are lanes downstream of certification, not inputs to
it, and their non-execution here is reported as an open item rather than
elided.

\paragraph{(c) External-actuation packets.}
\texttt{just status} in this sandbox reports \texttt{PUBLICATION\_ACTUATION=PENDING\_EXTERNAL\_ACTUATION},
with the arXiv submission packet (\texttt{ARXIV\_PACKET}) and the
GitHub-push packet (\texttt{GITHUB\_PUSH\_PACKET}) both currently
\texttt{BLOCKED}, while the GitHub-release packet
(\texttt{GITHUB\_RELEASE\_PACKET}) reports \texttt{ALIVE}. In this
repository's status taxonomy, \texttt{PENDING\_EXTERNAL\_ACTUATION}
denotes a packet that is itself complete but whose remaining step ---
submitting to arXiv, pushing to a remote the agent does not have
standing to push to --- is the user's action, not a further manufacturing
step this pipeline can discharge on its own. Reporting these packets as
blocked or pending is, consistent with this thesis's discipline
throughout, preferable to describing the release as ``published'' before
that external step has actually occurred.

\paragraph{(d) A sandbox toolchain limitation, scoped to this environment.}
This chapter closes with an honest report of a limitation particular to
the sandbox this thesis was written in, stated carefully so that it is
not read as a claim about the certified \ReleaseTag\ release itself.
Checked in this sandbox on 2026-07-21: \texttt{lake} and \texttt{elan}
are not on \texttt{PATH}, \texttt{ggen} is not on \texttt{PATH}, and
\texttt{procint}'s and \textsf{mfact}'s pinned toolchain
(\texttt{leanprover/lean4:v4.31.0}, per \texttt{lean-toolchain}) is not
installed in this sandbox's \texttt{elan} toolchain list, which is empty.
\texttt{just doctor}, itself a read-only diagnostic that does not dirty
the tree, confirms each of these and additionally reports that the
\texttt{lean-math-pack}, \texttt{quadrature-pack}, and
\texttt{post-release-pack} sources are absent, because they live at a
\texttt{/Users/sac/mfact} path this sandbox --- rooted at
\texttt{/home/user/mfact} --- does not have. A single bounded
installation attempt, \texttt{just install}, was run per this
repository's actuation discipline (Chapter~\ref{ch:ai-development}); it
proved to be a checker rather than an installer, printing guidance
(\texttt{elan: not found --- install with: curl https://elan.dev/install.sh -sSf | sh};
an analogous \texttt{ggen} instruction pointing at \texttt{just install-ggen}
from a praxis checkout) without itself performing a network fetch, and no
further installation was attempted, consistent with this repository's
instruction not to bypass its own guarded actuation paths. Toolchain
availability in this sandbox was unchanged, unavailable both before and
after this bounded attempt.

None of this bears on the certified release's own standing. \texttt{just status},
read-only in the same sandbox, reports the core release --- tag
\ReleaseTag, \TotalDecls\ declarations, \KernelTheorems\ proven,
\StatedCount\ stated, all gates \CertifiedStatus, tree clean, and a
93-artifact ledger --- as already certified from prior, checked-in
release state; nothing in this session rebuilt or re-derived that
standing, and nothing in this sandbox's missing toolchain retroactively
unproves a theorem the pinned toolchain already admitted at
certification time. STANDING.md is the governing record of that
certification, independent of any later environment's ability to replay
it locally. The limitation reported here is a property of this
particular writing environment's bootstrap state, worth recording
honestly precisely because this thesis's own discipline forbids asserting
replay success without having actually run it --- not a property, claimed
or implied, of the release the rest of this thesis reports on.

\chapter{Conclusion}
\label{ch:conclusion}

This thesis opened by separating \emph{contact} from \emph{consequence}:
an untrusted candidate producer --- human editor, template engine, or
large language model --- may touch a formal repository at many points,
but contact does not by itself confer standing on what it produces. Only
a declared, mechanically checked admission path can do that, and the
standing law $R_B \vdash A = \mu(O^*_B)$ names what such a path must
supply: a receipt~$R_B$ showing that an artifact~$A$ with standing is the
lawful output of a manufacturing transform~$\mu$ applied to admitted
observations~$O^*_B$, inside a declared boundary~$B$. The thesis argued
that this discipline is operational, not aspirational, by exhibiting it
across five pipeline stages --- source declaration in Turtle, ggen
projection, Lean kernel admission, \textsf{mfact} certification, and a
replayable receipt --- and by exercising the full pipeline on a
substantial first manufactured vertical, \textsf{procint}.

The evidentiary basis for that argument has been reported throughout
this thesis in the same way it is reported in \texttt{paper/main.tex}
itself: as rendered, ledgered fragments of the release manifest, never as
hand-typed numbers. \textsf{procint}'s catalog, at \TotalDecls\
declarations, reaches \KernelTheorems\ kernel-admitted, axiom-audited
theorems and \StatedCount\ formally stated declarations whose proof
obligations remain open, with the crown-jewel workflow-net soundness
equivalence carrying \CrownJewelStatus\ status under its
\texttt{[Finite~T]} hypothesis and a deliberately unbounded countermodel
witnessing why that hypothesis cannot be dropped. Chapter~\ref{ch:correspondence}
extended this discipline outward from Lean's own world of hand-written
types to a second rail, correspondence with Aeneas/Charon-extracted Rust,
governed by the same requirement that a correspondence theorem actually
quantify over its extracted declaration rather than a same-shaped
stand-in. Chapter~\ref{ch:empirical} extended it again, to praxis-graphlaw
and rslab's receipted empirical evidence, kept deliberately apart from
\textsf{procint}'s kernel-checked proven status by rslab's own
declared~$<$~extracted~$<$~stated~$<$~proven ladder. Chapter~\ref{ch:evaluation}
consolidated this evidence and closed with the eight enumerated failure
classes whose empty inhabitation, after certification, is itself checked
by negative controls rather than asserted. Chapter~\ref{ch:discussion}
then turned the same scrutiny back on the pipeline's own boundaries:
what proven-versus-stated does and does not claim, what kernel admission
does and does not check about statement fidelity, and an honest,
carefully scoped report of this particular writing environment's own
toolchain-bootstrap limitation, kept distinct from the certified
release's already-established standing.

Read together, the chapters support a narrower and more defensible claim
than ``AI-generated formal mathematics is trustworthy.'' They support the
claim that a manufacturing pipeline can make an untrusted producer's
output \emph{accountable}: every admitted theorem traces to a source
declaration, a projection template, a kernel-checked term, and a manifest
entry that a third party can re-check; every stated-but-unproven
declaration is labeled as such rather than allowed to read as proven;
every deliberate negative witness is guarded against promotion; and every
claim this thesis or its accompanying paper makes about counts, hashes,
or status is a citation of a receipted artifact rather than a fact
asserted on the strength of having been written down. Chapter~\ref{ch:discussion}'s
exclusions and residual risks --- most centrally, that kernel admission
does not by itself establish statement fidelity to an informal target ---
are not loose ends this conclusion resolves; they are the honest edge of
what this manufacturing discipline currently accomplishes, reported as
such rather than smoothed over, in keeping with the same discipline the
rest of this thesis has tried to exercise on every number it has cited.


\bibliographystyle{plain}
\bibliography{../paper/refs}

\end{document}
